% RDRS Bangladesh — MIS Internship Report (polished, locally + Overleaf compilable)
% Local: TeX Live base only (color/colortbl/fancyhdr). Overleaf: auto-enables hyperref/booktabs/microtype if present.
% Compile: tectonic RDRS_MIS_Internship_Report.tex
\documentclass[12pt,a4paper]{report}

\usepackage[utf8]{inputenc}
\usepackage[T1]{fontenc}
\usepackage{lmodern}
\usepackage[a4paper,margin=2.4cm,headheight=15pt,includehead,includefoot]{geometry}

% ---------- Base packages (guaranteed local) ----------
\usepackage{graphicx,color,colortbl,array,longtable,multicol,amsmath,amssymb,url}

% ---------- Colors (color.sty rgb, works everywhere) ----------
\definecolor{RDRSBlue}{rgb}{0.12,0.22,0.39}
\definecolor{RDRSTeal}{rgb}{0.05,0.49,0.48}
\definecolor{RDRSOrange}{rgb}{0.90,0.44,0.13}
\definecolor{RDRSLight}{rgb}{0.95,0.96,0.97}
\definecolor{RDRSGrey}{rgb}{0.35,0.35,0.35}
\definecolor{RDRSMidGrey}{rgb}{0.62,0.64,0.67}

% ---------- Optional enhancements (Overleaf/full TeX only, skipped locally) ----------
\IfFileExists{etoolbox.sty}{\usepackage[colorlinks=true,linkcolor=RDRSBlue,citecolor=RDRSTeal,urlcolor=RDRSBlue,breaklinks=true]{hyperref}}{}
\IfFileExists{booktabs.sty}{\usepackage{booktabs}}{%
  \newcommand{\toprule}{\hline\hline}%
  \newcommand{\midrule}{\hline}%
  \newcommand{\bottomrule}{\hline\hline}%
  \newcommand{\cmidrule}[1]{\hline}%
}
\IfFileExists{microtype.sty}{\usepackage{microtype}}{}
\IfFileExists{setspace.sty}{\usepackage{setspace}\setstretch{1.32}}{\renewcommand{\baselinestretch}{1.32}}

% ---------- Typography & page discipline ----------
\setlength{\parskip}{0.55em}
\setlength{\parindent}{0em}
\setlength{\emergencystretch}{3em}
\tolerance=1200
\clubpenalty=10000
\widowpenalty=10000
\raggedbottom
\setlength{\tabcolsep}{6pt}
\renewcommand{\arraystretch}{1.35}
\arrayrulecolor{RDRSMidGrey}
\setlength{\arrayrulewidth}{0.5pt}
\setlength{\abovecaptionskip}{6pt}
\setlength{\belowcaptionskip}{10pt}

% ---------- Headers / footers ----------
\usepackage{fancyhdr}
\pagestyle{fancy}
\fancyhf{}
\fancyhead[L]{\small\textsc{\color{RDRSBlue}RDRS Bangladesh} \color{RDRSMidGrey}\small$|$ \small MIS Internship Report}
\fancyhead[R]{\small\color{RDRSGrey}Mahmud (12022065)}
\fancyfoot[C]{\small\color{RDRSGrey}\thepage}
\renewcommand{\headrulewidth}{0.8pt}
\renewcommand{\footrulewidth}{0pt}
% color the head rule blue
\makeatletter
\def\headrule{{\if@fancyplain\let\headrulewidth\plainheadrulewidth\fi\color{RDRSBlue}\hrule\@height\headrulewidth\@width\headwidth \vskip-\headrulewidth}}
\makeatother

% ---------- Chapter / section styling (no titlesec needed) ----------
\makeatletter
\def\@makechapterhead#1{%
  \vspace*{10\p@}%
  {\parindent \z@ \raggedright \normalfont
    {\color{RDRSOrange}\bfseries\small\MakeUppercase{\@chapapp\ \thechapter}\par\nobreak}
    \vskip 4\p@
    {\color{RDRSBlue}\interlinepenalty\@M \Huge \bfseries #1\par\nobreak}%
    \vskip 6\p@
    {\color{RDRSBlue}\rule{\textwidth}{1.2pt}}\par
    \vskip 14\p@
  }}
\def\@makeschapterhead#1{%
  \vspace*{10\p@}%
  {\parindent \z@ \raggedright \normalfont
    {\color{RDRSBlue}\interlinepenalty\@M \Huge \bfseries #1\par\nobreak}%
    \vskip 6\p@
    {\color{RDRSBlue}\rule{\textwidth}{1.2pt}}\par
    \vskip 14\p@
  }}
\renewcommand{\section}{\@startsection{section}{1}{\z@}{-2.8ex plus -0.8ex minus -.2ex}{1.4ex plus .2ex}{\normalfont\Large\bfseries\color{RDRSBlue}}}
\renewcommand{\subsection}{\@startsection{subsection}{2}{\z@}{-2.2ex plus -.6ex minus -.2ex}{1.0ex plus .2ex}{\normalfont\large\bfseries\color{RDRSBlue}}}
\renewcommand{\subsubsection}{\@startsection{subsubsection}{3}{\z@}{-1.6ex plus -.4ex minus -.2ex}{0.8ex plus .2ex}{\normalfont\normalsize\bfseries\color{RDRSTeal}}}
% captions: small, bold colored label
\long\def\@makecaption#1#2{%
  \vskip\abovecaptionskip
  \sbox\@tempboxa{\small\textbf{\color{RDRSBlue}#1.} \small\color{RDRSGrey}#2}%
  \ifdim \wd\@tempboxa >\hsize
    {\small\textbf{\color{RDRSBlue}#1.} \small\color{RDRSGrey}#2\par}%
  \else
    \global \@minipagefalse
    \hb@xt@\hsize{\hfil\box\@tempboxa\hfil}%
  \fi
  \vskip\belowcaptionskip}
\makeatother

% ---------- Reusable visual elements & callout boxes ----------
\setlength{\fboxsep}{10pt}
\setlength{\fboxrule}{0.9pt}
\newcommand{\highlightbox}[1]{\begin{center}\fcolorbox{RDRSBlue}{RDRSLight}{\begin{minipage}{0.92\textwidth}#1\end{minipage}}\end{center}}
\newcommand{\keybox}[1]{\begin{center}\fcolorbox{RDRSTeal}{RDRSLight}{\begin{minipage}{0.92\textwidth}#1\end{minipage}}\end{center}}
\newcommand{\warnbox}[1]{\begin{center}\fcolorbox{RDRSOrange}{RDRSLight}{\begin{minipage}{0.92\textwidth}#1\end{minipage}}\end{center}}
\newcommand{\headercell}[1]{\textcolor{white}{\textbf{#1}}}

% Quick chapter summary strip for skimmers
\newcommand{\chapterglance}[2]{%
  \begin{center}%
  \fcolorbox{RDRSBlue}{RDRSLight}{%
    \begin{minipage}{0.94\textwidth}%
      \textbf{\color{RDRSBlue}\small CHAPTER AT A GLANCE \quad $|$ \quad #1}\\[3pt]%
      \footnotesize\color{RDRSGrey}#2%
    \end{minipage}%
  }%
  \end{center}%
  \vspace{0.25cm}%
}

% Visual formula callout card
\newcommand{\formulacard}[4][RDRSBlue]{%
  \begin{center}%
  \fcolorbox{#1}{RDRSLight}{%
    \begin{minipage}{0.94\textwidth}%
      \centering%
      \small\textbf{\color{#1}#2}\\[4pt]%
      #3\\[4pt]%
      \footnotesize\color{RDRSGrey}#4%
    \end{minipage}%
  }%
  \end{center}%
}

% Visual strength / gap tags
\newcommand{\strengthtag}{\textcolor{RDRSTeal}{\textbf{[STRENGTH]}}}
\newcommand{\gaptag}{\textcolor{RDRSOrange}{\textbf{[GAP]}}}

\title{\textbf{Evaluating MIS Practices in Project Monitoring:}\\ \textbf{A Study on RDRS Bangladesh}}
\author{A.~R.~Mahmud\\ID: 12022065\\BBA, Department of Management Information Systems\\Begum Rokeya University, Rangpur}
\date{}

\begin{document}

% ================= TITLE PAGE =================
\begin{titlepage}
\centering
\vspace*{0.3cm}
{\small\scshape\color{RDRSTeal} Begum Rokeya University, Rangpur \quad $\bullet$ \quad Department of Management Information Systems\par}
\vspace{0.35cm}
{\color{RDRSOrange}\rule{\textwidth}{2.2pt}}\\[2pt]
{\color{RDRSBlue}\rule{\textwidth}{0.7pt}}
\vspace{0.5cm}
{\small\bfseries\color{RDRSOrange} \MakeUppercase{Internship Report}\par}
\vspace{0.15cm}
{\LARGE\bfseries\color{RDRSBlue} Evaluating MIS Practices\\ in Project Monitoring\par}
\vspace{0.15cm}
{\Large\color{RDRSGrey} A Study on RDRS Bangladesh\par}
\vspace{0.3cm}
{\color{RDRSBlue}\rule{3cm}{2.2pt}}
\vspace{0.55cm}

\fcolorbox{RDRSBlue}{RDRSLight}{%
\begin{minipage}{0.90\textwidth}
\centering
\small\scshape\color{RDRSBlue} Submitted By \quad $/$ \quad Submitted To \quad $/$ \quad Host\par
\vspace{0.15cm}
\begin{tabular}{ccc}
\parbox{4.6cm}{\centering \textbf{A.~R.~Mahmud}\\ ID: 12022065\\ \small BBA, Major in MIS} &
\parbox{0.6cm}{\centering \color{RDRSMidGrey}$|$} &
\parbox{5.2cm}{\centering \textbf{Mohammad Rakibul Islam Bhuiyan}\\ \small Assistant Professor, MIS, BRUR} \\
\end{tabular}\\[6pt]
{\small\color{RDRSGrey} In collaboration with \textbf{\color{RDRSBlue}RDRS Bangladesh, MEL Department}\par}
\end{minipage}}
\vspace{0.45cm}

\fcolorbox{RDRSTeal}{RDRSTeal}{\begin{minipage}{0.90\textwidth}\centering\textcolor{white}{\small\bfseries Internship at RDRS Bangladesh \quad $\bullet$ \quad Fieldwork: Rangpur \quad $\bullet$ \quad September--November 2026}\end{minipage}}

\vfill
{\small\color{RDRSGrey} Date of Submission: 30 November 2026 \qquad Place: Rangpur \par}
\vspace{0.15cm}
{\scriptsize\color{RDRSMidGrey} Sources: MEL monitoring schedule \quad $\bullet$ \quad G~Banker outputs \quad $\bullet$ \quad C-RANGPUR instrument \& F01-112 matrix \quad $\bullet$ \quad KII \& field notes \quad $\bullet$ \quad RDRS Strategic Plan \quad $\bullet$ \quad Somriddhi report \quad $\bullet$ \quad Court audit \par}
\end{titlepage}

% ================= FRONT MATTER =================
\pagenumbering{roman}

\chapter*{Letter of Transmittal}
\addcontentsline{toc}{chapter}{Letter of Transmittal}
\enlargethispage{2\baselineskip}
30 November 2026\par
Mohammad Rakibul Islam Bhuiyan\\
Assistant Professor, Department of Management Information Systems\\
Begum Rokeya University, Rangpur\\[4pt]
Subject: Submission of Internship Report on ``Evaluating MIS Practices in Project Monitoring: A Study on RDRS Bangladesh''\\[4pt]
Dear Sir,\\[2pt]
I am pleased to submit my internship report, prepared as a mandatory requirement for completing the BBA program in Management Information Systems at Begum Rokeya University, Rangpur. The report examines the practical MIS and monitoring workflows of RDRS Bangladesh, based on my three-month placement (01 September to 30 November 2026) within the Monitoring, Evaluation, and Learning (MEL) Department at the Central Coordination Office in Rangpur.

Throughout this placement, I observed how field data travels from remote village groups and branch offices to central management dashboards. I had the opportunity to study the C-RANGPUR baseline survey instrument, live G~Banker microfinance reports, the ICPS child protection baseline, the Access to Justice court tracking system, and internal monitoring archives. This report summarizes the operational strengths, data integrity gaps, and practical lessons I gathered from staff interviews, document reviews, and field observations.

I have made every effort to keep this analysis objective, evidence-based, and grounded in the actual documents I inspected. I am sincerely grateful for your supervision, constructive feedback, and continuous encouragement during this study.\\[6pt]
Sincerely,\\
\textbf{A.~R.~Mahmud} \quad (ID: 12022065)\\
Department of Management Information Systems\\
Begum Rokeya University, Rangpur

\chapter*{Acknowledgement}
\addcontentsline{toc}{chapter}{Acknowledgement}
\enlargethispage{2\baselineskip}
I would like to express my sincere gratitude to my academic supervisor, Mohammad Rakibul Islam Bhuiyan, Assistant Professor, Department of Management Information Systems, Begum Rokeya University, Rangpur, for his guidance, structural advice, and review throughout the preparation of this report.

I am deeply indebted to the management and staff of RDRS Bangladesh at the Central Coordination Office, Rangpur, for hosting my three-month placement and providing generous access to operational documents and systems:
\begin{itemize}
\setlength{\itemsep}{1pt}\setlength{\parskip}{0pt}
\item \textbf{Mr.~Mamunur Rashid} (Head of MEL) and \textbf{Dabasish Mahata} (Coordinator-Monitoring) for approving my placement schedule, providing departmental orientation, and supervising my weekly learning objectives.
\item \textbf{Md.~Saiful Alam (Shekul)} (ICPS Baseline Team; ex-MEAL \& GIS, Save the Children) for a detailed walkthrough of baseline indicator design, ESS safeguards, and MEL maturity.
\item \textbf{Fahmida Mostafa} (Senior MEL Officer) for explaining active project portfolios, routine field reporting workflows, and institutional scale across the five core program domains.
\item \textbf{Mr.~Iftekharul Islam} (Manager-Monitoring) for sharing the thirty-year historical evolution of microfinance software at RDRS, from early manual ledgers to Micro Bank and modern G~Banker automation.
\item \textbf{Md.~Ajazul Haque} (Manager, MEL) for walking me through the C-RANGPUR digital survey pipeline, KoBo Toolbox architecture, and Household Climate Vulnerability Index (HCVI) scoring formulas.
\item \textbf{Moushumi Akter}, \textbf{Maydul Islam}, \textbf{Md.~Ziaur Rahman}, \textbf{Dr.~Salah Uddin}, \textbf{Rashedul Arefeen}, \textbf{Rumana Azim Ansari}, and \textbf{Maj Shammita Kabir (Retd)} for their time and guidance during field visits and departmental rotations.
\end{itemize}
Finally, I thank my family and classmates for their constant encouragement. Any errors in interpretation or analysis in this report remain solely mine.

\chapter*{Executive Summary}
\addcontentsline{toc}{chapter}{Executive Summary}
\highlightbox{\small\textbf{\color{RDRSBlue}Guiding Question:} How does monitoring data actually move from grassroots village groups to executive management decisions at RDRS Bangladesh, and where does the information chain hold or break?}

This report documents what I observed during my three-month internship (01 September to 30 November 2026) based at the RDRS Central Coordination Office in Rangpur. Rather than relying on promotional brochures, I tracked the actual tools field staff carry: physical registers and structured Excel workbooks for weekly branch operations, KoBo Toolbox on mobile phones for ad-hoc household surveys, and central dashboards built in Power~BI and G~Banker for management review. My placement followed an official rotation sanctioned by the Head of MEL (reproduced in Appendix~\ref{sec:official-schedule}), covering baseline surveys, branch audits, and cross-departmental operations across HR, Agriculture, Social Development, Safeguarding, and General Administration.

The findings reveal two contrasting data realities within the organization:

First, microfinance operations have matured into a centralized, automated pipeline. Between 1995 and 2011, RDRS relied heavily on a legacy software system called Micro Bank, which created substantial organizational risk because reporting depended almost entirely on one or two senior individuals. The transition to South Tech in 2011 and subsequently to G~Banker in 2017/18 eliminated this vulnerability. Today, branch-level balance sheets (POMIS-5), income statements (POMIS-4), and Credit Organizer (CO) performance reports generate with a single click. The key operational lesson here is institutional continuity: software works best when indicator rules (such as Portfolio-at-Risk and On-Time Repayment) are codified and knowledge is shared across multiple officers.

Second, development projects operate on a much more fragmented stack. In projects like C-RANGPUR, field enumerators collect rich data through a 10-section KoBo digital questionnaire, calculating complex indicators and a composite Household Climate Vulnerability Index (HCVI, weighted 30\% Exposure, 30\% Sensitivity, and 40\% Adaptive Capacity Deficit). However, data processing is divided across three separate environments---KoBo servers, manual Excel clean-up, and Power~BI visualizations---with no automated connection to the core microfinance database.

The most critical vulnerabilities appear at the ``last mile'' of field reporting. A review of internal monitoring reports (such as the May 2023 Somriddhi audit) showed persistent mismatches between physical paper registers and digital central reports (e.g., 89\% actual student attendance in field registers entered as 100\% in the central MIS; missing expenditure files for 65\% of livelihood asset grants). High staff turnover and paper-heavy workflows in Credit Shield microinsurance further slow down operations.

\keybox{\small\textbf{\color{RDRSTeal} Priority Recommendations:}
\textbf{(1) Mandatory Register-to-MIS Reconciliation Sign-off:} Require branch supervisors to physically verify and cross-sign field registers against digital entries before monthly reports lock.
\textbf{(2) End-to-End KoBo Pipeline:} Automate indicator scoring directly from mobile survey submissions into Power~BI to eliminate error-prone intermediate Excel copies.
\textbf{(3) Documented Indicator Ownership:} Maintain formal standard operating procedures and dual ownership for every recurring report family to prevent single-person dependencies.}

\tableofcontents
\listoffigures
\listoftables

\highlightbox{\centering\small \textbf{\color{RDRSBlue}Abbreviations:} MEL -- Monitoring, Evaluation and Learning; MEAL -- plus Accountability; ESS -- Environmental and Social Safeguard; TOR -- Terms of Reference; FGD -- Focus Group Discussion; KII -- Key Informant Interview; HH -- Household; IGA -- Income Generating Activity; CSA -- Climate Smart Agriculture; HCVI -- Household Climate Vulnerability Index (EI -- Exposure Index; SI -- Sensitivity Index; ACI -- Adaptive Capacity Index); NCD -- Non-Communicable Disease; SPSS -- Statistical Package for the Social Sciences; PKSF -- Palli Karma-Sahayak Foundation; CO -- Credit Organizer; OTR -- On-Time Repayment; PAR -- Portfolio at Risk; DQA -- Data Quality Assessment; ATM -- Assistant Thana Manager (historic); MSDS -- Monthly Savings Deposit Scheme; WASH -- Water, Sanitation and Hygiene; LRL -- Livelihood Restoration Loan; BFIU -- Bangladesh Financial Intelligence Unit; A2J -- Access to Justice for Women; GIZ -- Deutsche Gesellschaft f{\"u}r Internationale Zusammenarbeit; ADR -- Alternative Dispute Resolution; DLAO -- District Legal Aid Officer / Office; DLAC / UZLAC / UPLAC -- District / Upazila / Union Parishad Legal Aid Committee; LSID -- Legal Service Information Desk; PNGO -- Partner Non-Governmental Organization; CDR -- Community Dispute Resolution; GRO -- General Register Officer (Court Police); ATT / CT / AHTT / ST -- Anti-Terrorism / Cyber / Anti-Human Trafficking / Special Tribunal; CrPC -- Code of Criminal Procedure; DBLA -- District Bar Legal Aid.}

\clearpage
\pagenumbering{arabic}

% ================= CHAPTER 1 =================
\chapter{Introduction}

\chapterglance{Background, Objectives \& Methodology}{Evaluates how monitoring data is collected, verified, and utilized across RDRS Bangladesh projects. Grounded in a structured 3-month field placement (01 Sept -- 30 Nov 2026) following the official MEL internship schedule (Appendix~\ref{sec:official-schedule}).}

\section{Background of the Study}
When studying Management Information Systems in a BBA classroom, textbooks describe seamless databases where transactions flow instantly into executive dashboards. Real-world development work in rural Bangladesh is far messier. Systems must survive intermittent mobile networks, handwritten field registers, tight donor deadlines, and staff turnover across hundreds of remote branch offices.

I chose to intern with the Monitoring, Evaluation, and Learning (MEL) Department of RDRS Bangladesh to observe how a major non-governmental organization actually monitors field progress. RDRS operates across 33 districts and reaches over 1.6 million people. Tracking interventions across that scale requires a delicate balance between field registers, mobile data collection, and central management dashboards. This report investigates that information plumbing: who records the numbers, which forms or apps they use, how data gets validated, and where reporting breakdowns occur before reaching executive decision-makers.

\section{Objectives of the Study}
\textbf{Broad Objective:} To evaluate the practical management information system (MIS) practices, tools, and workflows used in project monitoring at RDRS Bangladesh.

\textbf{Specific Objectives:}
\begin{enumerate}
\item To identify the specific digital and manual MIS tools used by the MEL Department.
\item To map how monitoring data travels from grassroots field units to central management dashboards.
\item To examine operational strengths, data integrity gaps, and structural bottlenecks in current monitoring practices.
\item To formulate practical, field-grounded recommendations for strengthening data reliability and institutional continuity.
\end{enumerate}

\section{Scope of the Study}
The study focuses on the practices observed during my three-month placement (01 September to 30 November 2026) at the RDRS Central Coordination Office in Rangpur. I had direct access to files, schedules, and personnel connected to several active initiatives: the C-RANGPUR climate resilience baseline survey, the ICPS child protection project, maternal health monitoring under Jononni, adolescent education under SEEDS, the formal justice audit under Access to Justice for Women (A2J), the core microfinance portfolio (Udayan and Unnoti programs), and archived monitoring reports from the PKSF-funded Somriddhi project. The study does not assess external financial audits, statutory governance reviews, or projects outside the MEL Department's active monitoring mandate.

\section{Methodology}
\subsection{Placement structure and field exposure}
My internship followed a formal three-month training plan sanctioned by Mr.~Mamunur Rashid, Head of MEL (see Appendix~\ref{sec:official-schedule} for the complete 13-stage program). The placement was organized into six distinct phases:
\begin{enumerate}
\item \textbf{Orientation and Guidelines (01--10 Sept 2026):} Overview of RDRS organizational history, five core program domains, and departmental reporting procedures with Mr.~Mamunur Rashid, Dabasish Mahata, Mr.~Iftekharul Islam, and Fahmida Mostafa.
\item \textbf{Baseline Fieldwork Rotation (13--17 Sept 2026):} Field exposure to baseline data collection under the C-RANGPUR project with Moushumi Akter (Monitoring and Documentation Officer).
\item \textbf{Branch Monitoring and Verification (20--24 Sept 2026):} Field inspection in the Rangpur Region observing passbook audits, ledger entries, and manual-to-digital loan reconciliation with Maydul Islam (Senior Monitoring Officer) and Md.~Ziaur Rahman (Senior Officer, Admin \& MIS).
\item \textbf{Cross-Departmental Institutional Rotations (27 Sept -- 01 Oct 2026):} Targeted briefings across Human Resources (Md.~Ziaur Rahman), Agriculture and Climate Change / Rupantar Project (Dr.~Salah Uddin and Pallab Goswami), Social Development and Union Federations (Rashedul Arefeen), Safeguarding Policy (Rumana Azim Ansari), and Admin / General Services (Maj Shammita Kabir, Retd).
\item \textbf{Deep MIS Practice in MEL (04 Oct -- 12 Nov 2026):} Hands-on examination of mobile survey design in KoBo Toolbox, indicator formula matrices, the Household Climate Vulnerability Index (HCVI), G~Banker microfinance exports, and court case tracking files with Dabasish Mahata and Md.~Ajazul Haque.
\item \textbf{Synthesis, Presentation, and Reporting (15--30 Nov 2026):} Report compilation, faculty consultation, internal learning presentation, and final submission to the Head of MEL.
\end{enumerate}

\subsection{Primary data sources}
Primary evidence came from structured walkthroughs, daily observational logs, and key informant interviews with the RDRS staff members who build, operate, and maintain these systems (Table~\ref{tab:interviews}).

\begin{table}[htbp]
\centering\small\caption{Key Informant Interviews Conducted}\label{tab:interviews}
\begin{tabular}{|>{\raggedright\arraybackslash}p{3.2cm}|>{\raggedright\arraybackslash}p{3.4cm}|>{\raggedright\arraybackslash}p{7.0cm}|}
\hline
\rowcolor{RDRSBlue}\headercell{Interviewee} & \headercell{Designation} & \headercell{Key Focus Area \& System Exposure} \\ \hline
Md.~Saiful Alam (Shekul) & ICPS Baseline Team (ex-MEAL \& GIS, Save the Children) & MEL functions; 6-step baseline design lifecycle; 21 UNICEF child protection indicators; ESS safeguards. \\ \hline
\rowcolor{RDRSLight} Fahmida Mostafa & Senior MEL Officer & Active project portfolio; departmental reporting workflow; 5 thematic areas; field scale and coverage. \\ \hline
Mr.~Iftekharul Islam & Manager, Monitoring & Historical evolution of microfinance software: manual ATM era $\to$ Micro Bank $\to$ South Tech $\to$ G~Banker. \\ \hline
\rowcolor{RDRSLight} Md.~Ajazul Haque & Manager (MEL) & C-RANGPUR digital pipeline; KoBo questionnaire logic; HCVI 30/30/40 scoring; Power~BI dashboard integration. \\ \hline
Project Focal \& Paralegals & Access to Justice for Women & Formal court system monitoring; physical docket enumeration across \textit{Ejlash} and \textit{Serestha}; court-annexed ADR. \\ \hline
\end{tabular}
\end{table}

\subsection{Secondary data sources}
Secondary materials reviewed during the placement included:
\begin{itemize}
\item The MEL Department September 2026 Monthly Monitoring Schedule (branch and project allocations).
\item Live G~Banker system exports: CO-wise Program Changes Report (Udayan office), POMIS-4 (Income/Expenditure Statement), and POMIS-5 (Statement of Financial Position).
\item C-RANGPUR baseline survey instrument (\url{https://ee-eu.kobotoolbox.org/x/wzjSgzZ6} and internal questionnaire specification document).
\item Project monitoring matrix spreadsheet file \texttt{F01-112} (tab ``Monitoring Planning'').
\item Internal policy manuals: RDRS 6th Strategic Plan (2023--2027), Annual Report 2021--22, Credit Shield Insurance Agreement (Guardian Life, July 2024), and the AML/CFT Compliance Policy (January 2022).
\item Internal monitoring archive: Somriddhi Programme Field Monitoring Report (May 2023, PKSF funding, covering Panchagarh and Kurigram).
\end{itemize}

\section{Limitations of the Study}
While I had open access within the MEL Department, several practical constraints shaped this study. First, the evaluation was conducted during an active academic internship term, limiting total field immersion to three months. Second, direct database backend access (SQL queries or raw server configuration for G~Banker and KoBo) was restricted to IT administrators for security reasons; my analysis is based on exported operational reports, front-end user forms, and system documentation. Third, ongoing baseline surveys (such as ICPS and C-RANGPUR) were still in their preliminary stages during the placement, meaning baseline values were being established rather than compared to final project outcomes. All figures and policy thresholds are explicitly cited with their source document date.

% ================= CHAPTER 2 =================
\chapter{Organizational Overview}

\chapterglance{Institutional Scale \& Operational Structure}{Profiles RDRS Bangladesh: 50-year evolution from post-war relief to 406 Union Federations, 33 districts, 1.64M participants, and BDT 15.09B in active microfinance.}

\section{About RDRS Bangladesh}
RDRS Bangladesh traces its roots directly to the aftermath of Bangladesh's Liberation War. In February 1972, the Lutheran World Federation (Department for World Service, Geneva) launched a field program in Rangpur to provide emergency relief, medical supplies, housing, and agricultural tools to war-affected communities. Originally called the Rangpur Dinajpur Relief Service, the organization subsequently evolved into the Rangpur Dinajpur Rural Service, reflecting a shift toward long-term rural development. In 1997, it registered as an independent Bangladeshi NGO under the Trust Act, governed by its own national Board of Trustees.

Geographically, RDRS concentrated its efforts in the northwestern districts of Bangladesh. This region has historically suffered from structural isolation, recurring floods from the Teesta--Jamuna--Brahmaputra river systems, severe riverbank erosion, and the seasonal lean period known as \textit{Monga}. Over five decades, the organization's work transitioned through distinct phases: post-war relief (1972--1975); health clinics, non-formal schools, and agricultural support (1976--1989); localization and federation building (1990--1997); and integrated microfinance, climate adaptation, and community empowerment under successive strategic plans. Celebrating its 50th anniversary in February 2022, RDRS is currently executing its 6th Strategic Plan (2023--2027). The organization maintains its central executive leadership at House 43, Road 10, Sector 6, Uttara, Dhaka, while operational field management is coordinated from its Central Coordination Office on Jail Road in Rangpur.

\begin{figure}[htbp]
\centering
\fcolorbox{RDRSBlue}{RDRSLight}{%
\begin{minipage}{0.92\textwidth}
\centering
\small\textbf{\color{RDRSBlue}1972: LWF Emergency Relief} \quad \textcolor{RDRSOrange}{$\blacktriangleright$} \quad \small\textbf{\color{RDRSBlue}1997: Trust Act Localization} \quad \textcolor{RDRSOrange}{$\blacktriangleright$} \quad \small\textbf{\color{RDRSBlue}2017/18: Automated G~Banker Era} \\[3pt]
{\scriptsize\color{RDRSGrey} Emergency War Relief $\longrightarrow$ Independent National NGO $\longrightarrow$ Integrated Digital Dashboards\par}
\end{minipage}}
\caption{Historical evolution of RDRS Bangladesh across fifty years of operations.}
\label{fig:timeline}
\end{figure}

\section{Mission, Vision, and Core Values}
\textbf{Vision:} A just and peaceful society where citizens live in dignity, free from poverty, distress, and ignorance in a sustainable environment (RDRS Strategic Plan 2023--2027).

\textbf{Mission:} RDRS works with disadvantaged and marginalized communities and their organizations to empower them, build resilience, and ensure sustainable improvements in their lives with dignity. In practice, this means directing credit, agricultural advice, and social support toward rural poor households---especially women---while helping local community institutions stand on their own without perpetual NGO aid.

\textbf{Core Values:} The organization operates under five stated corporate values: \textit{Compassion, Commitment, Innovation, Adaptability, and Accountability}. During my placement, I observed that field staff often summarize these values through practical field rules: strict gender equity, transparent loan ledgers, and responsive grievance handling. For an MIS student, these values are directly relevant: every monitoring form, ledger entry, and dashboard is essentially a mechanism for institutional accountability.

\section{Field Scale and Geographic Coverage}
The operational scale of RDRS Bangladesh is extensive. Table~\ref{tab:scale-summary} summarizes the footprint data provided by the MEL Department and verified against recent institutional records.

\begin{table}[htbp]
\centering\small\caption{Institutional Footprint and Program Reach}\label{tab:scale-summary}
\begin{tabular}{|>{\raggedright\arraybackslash}p{4.2cm}|>{\raggedright\arraybackslash}p{3.8cm}|>{\raggedright\arraybackslash}p{6.4cm}|}
\hline
\rowcolor{RDRSBlue}\headercell{Operational Dimension} & \headercell{Scale / Quantity} & \headercell{Coverage Details \& Field Locations} \\ \hline
Union Federations & 406 Federations & Community-owned apex bodies at Union level \\ \hline
\rowcolor{RDRSLight} Administrative Reach & 33 Districts, 183 Upazilas & 872 Unions across Bangladesh \\ \hline
Program Participants & 1,642,724 Individuals & $\approx$1.64 million direct participants reached \\ \hline
\rowcolor{RDRSLight} Core Operations Zone & Rangpur Division & Kurigram, Gaibandha, Thakurgaon, Dinajpur, Nilphamari, Panchagarh, Rangpur, Lalmonirhat \\ \hline
Specialized Field Areas & Humanitarian \& Regional & Ukhiya and Teknaf (Cox's Bazar); targeted pockets in Sylhet, Mymensingh, and Rajshahi \\ \hline
\rowcolor{RDRSLight} Microfinance Network & 288 Branch Offices & 29 districts; 21,134 groups; 458,819 members; 377,250 active borrowers \\ \hline
Microfinance Portfolio & Financial Metrics & BDT~5.14B member savings; BDT~15.09B outstanding loans; 95.08\% recovery rate \\ \hline
\end{tabular}
\end{table}

\section{Governance and Departmental Organization}
Governance is provided by an independent Board of Trustees and General Body, with executive execution headed by the Executive Director. Foreign-funded development projects comply with statutory oversight under the Foreign Donations Regulation Act 2016 through the NGO Affairs Bureau (FD-6 approvals and mandatory annual audits). Microcredit operations are licensed and regulated by the Microcredit Regulatory Authority (MRA).

The field operational pyramid flows through three tiers: Regions (16--17 regional offices), Areas (59--61 area offices), and Branches (278--288 branch offices). Union Federations operate at the grassroots Union Parishad level as autonomous, member-run organizations that negotiate government safety-net claims, conduct local mediation, and hold seats on local standing committees.

\section{The MEL Department}
The Monitoring, Evaluation, and Learning (MEL) Department serves as the central data and quality assurance engine of RDRS. Led by the Head of MEL at the Rangpur Central Coordination Office, the department operates independently of direct program implementation teams to preserve objectivity.

Monitoring Officers receive regular reporting feeds from branch managers, credit organizers, project field officers, and community federations. Their core responsibility is straightforward: verify field registers against submitted summaries, conduct field spot-checks, compile divisional datasets, and produce monthly analytical reports for senior management. If a branch enters an erroneous figure or an enumerator skips a question in a survey, it is the MEL Department's job to catch the anomaly before reports reach the Executive Director.

\section{The Five Core Program Domains}
RDRS organizes all operational activities across five strategic thematic domains:
\begin{enumerate}
\item \textbf{Skills, Employment, and Entrepreneurship:} Vocational trades, market linkages, youth skills training, and enterprise development.
\item \textbf{Social Empowerment:} Union Federation strengthening, legal aid, gender justice, anti-child marriage advocacy, and civic participation.
\item \textbf{Agriculture, Food Security, and Climate Change:} Flood-resilient farming, crop diversification, renewable energy, and community adaptation.
\item \textbf{Microfinance and Employment Creation:} Collateral-free group loans, micro-enterprise credit, voluntary savings schemes, and microinsurance.
\item \textbf{Emergency and Humanitarian Response:} Disaster relief, post-flood rehabilitation in northern chars, and refugee support in Cox's Bazar.
\end{enumerate}

\begin{table}[htbp]
\centering\small\caption{Key Active Projects Observed During Placement}\label{tab:projects}
\begin{tabular}{|>{\raggedright\arraybackslash}p{2.2cm}|>{\raggedright\arraybackslash}p{4.8cm}|>{\raggedright\arraybackslash}p{7.4cm}|}
\hline
\rowcolor{RDRSBlue}\headercell{Project Name} & \headercell{Operational Focus} & \headercell{Donor Partner, Duration \& Scale} \\ \hline
Jononni & Maternal and neonatal health; safe delivery care; pregnant women tracking & Save the Children; March 2023--December 2027; active in Rangpur and Lalmonirhat \\ \hline
\rowcolor{RDRSLight} SEEDS & Socio-economic empowerment for adibashi adolescent girls; schooling and livelihoods & Str{\o}mme Foundation; January 2024--December 2028; 3,000 target households \\ \hline
MCGL & Preventable maternal and child mortality; public health systems strengthening & U.S. Department of State via Jhpiego; active in Panchagarh and Rangpur \\ \hline
\rowcolor{RDRSLight} ICPS & Disability-inclusive child protection; reducing violence against children with disabilities & UNICEF; June 2026--December 2027; Budget: BDT~2,20,40,834 (see Section~\ref{sec:icps}) \\ \hline
C-RANGPUR & Climate resilience and poverty upliftment; CSA adoption and household vulnerability & RDRS MEL Department; baseline launched September 2026 (see Section~\ref{sec:crangpur-instrument}) \\ \hline
\rowcolor{RDRSLight} A2J for Women & Access to Justice; formal court legal aid, court-annexed mediation, and docket audits & Law \& Justice Division (MoLJPA) \& German GIZ; 1 Mar 2024--31 Dec 2026 (extended to Mar 2027); BDT~1.82M (Section~\ref{sec:a2j}) \\ \hline
\end{tabular}
\end{table}

\section{Microfinance Products: Udayan, Unnoti, and Savings Schemes}
Understanding RDRS's monitoring systems requires understanding its microcredit structure. The organization does not issue a single generic loan; it manages specialized financial products designed for different stages of borrower graduation:

\textbf{Udayan (Group-based Microcredit):} Group-liability lending designed for newly enrolled female borrowers living in rural villages. Samity members meet weekly, cross-guarantee each other's installments, and deposit mandatory weekly savings.

\textbf{Unnoti (Individual Enterprise Credit):} Higher-ceiling loans provided to graduated borrowers with proven repayment histories. Unlike Udayan, Unnoti loans are underwritten individually based on household business cash flow rather than joint group liability.

\textbf{Savings Architecture:} Financial discipline is anchored by compulsory member savings. Every member contributes at least BDT~50 weekly, with a mandatory monthly minimum of BDT~200 for Udayan borrowers and BDT~500 for Unnoti clients. Compulsory deposits earn a 6\% annual profit share. In addition, members can open General Voluntary Savings (flexible deposits up to BDT~10,000 earning 7\% simple interest) or the Monthly Savings Deposit Scheme (MSDS), which offers 7\% on 3-year term deposits and 7.5\% on 5-year maturities.

\textbf{Targeted Credit Lines:}
\begin{itemize}
\item \textbf{Water Credit:} A WASH-focused loan ranging from BDT~10,000 to 100,000, operating in 93 branches across Dinajpur, Nilphamari, Thakurgaon, Panchagarh, Habiganj, and Moulvibazar. The loan is tied directly to technical assistance for installing deep tube-wells and sanitary latrines.
\item \textbf{Livelihood Restoration Loan (LRL Phase 2):} A post-COVID recovery credit line of up to BDT~50,000 with a 6-month grace period, 24-month maximum tenure, and an 8\% declining-balance charge, repaid in 18 monthly installments to recapitalize vulnerable micro-enterprises.
\end{itemize}

\highlightbox{\small\textbf{\color{RDRSBlue}Why product differentiation matters for MIS:} A group-based Udayan loan requires tracking weekly samity collections and cross-guarantee ledgers. An individual Unnoti loan requires monitoring cash-flow solvency and individual collateral. Water Credit requires verifying physical sanitary installations before funds disburse. The field monitoring checklist analyzed in Chapter~\ref{chap:analysis} is designed around these operational distinctions.}

% ================= CHAPTER 3 =================
\chapter{Literature Review}

\chapterglance{Theoretical Framework \& Field Realities}{Synthesizes MIS definitions, OECD-DAC monitoring standards, and empirical findings on microfinance software continuity and data integrity.}

\section{Management Information Systems in Practice}
In academic literature, Laudon and Laudon (2020) define a Management Information System as an integrated set of components---hardware, software, human resources, and business processes---that collect, store, process, and distribute data to support organizational decision-making and managerial control. In commercial corporate environments, MIS discussions frequently center on enterprise resource planning (ERP) architectures, cloud computing, and real-time database transactions.

In the non-governmental and microfinance sectors, however, the practical challenge of MIS is fundamentally different. As observed during my placement, the primary challenge is not procuring sophisticated software; it is ensuring the integrity and timeliness of data entered at the grass roots. A central dashboard is only as reliable as the weekly collection sheet filled out by a Credit Organizer in a remote village or the survey submitted by a field enumerator on a smartphone. For development managers, MIS is less about raw computational power and more about verification discipline, standard operational procedures, and data quality controls that prevent corrupted numbers from corrupting decisions.

\section{Project Monitoring Frameworks}
The Organisation for Economic Co-operation and Development's Development Assistance Committee (OECD-DAC) defines project monitoring as a continuous, systematic process that tracks the implementation progress of an intervention against predefined indicator targets. Unlike periodic evaluation, which measures long-term impact and sustainability at set intervals, monitoring is an active management tool designed to identify operational deviations, resource bottlenecks, and implementation shortfalls while activities are still underway.

A robust monitoring system relies on several core components:
\begin{itemize}
\item \textbf{The Baseline Study:} Establishes the empirical reference value for every project indicator before interventions begin.
\item \textbf{Data Quality Assessments (DQA):} Structured protocols that evaluate data across five standard dimensions: validity, reliability, integrity, precision, and timeliness.
\item \textbf{Mixed-Method Data Capture:} Combining quantitative household sample surveys (measuring statistical coverage and adoption rates) with qualitative Focus Group Discussions (FGD) and Key Informant Interviews (KII) to capture contextual nuances and beneficiary feedback.
\end{itemize}

\section{MIS Lessons from NGO and Microfinance Environments}
Empirical studies of information systems within microfinance institutions (MFIs) and international NGOs highlight three recurring vulnerabilities that mirror my observations at RDRS Bangladesh:

\begin{enumerate}
\item \textbf{The Risk of Person-Dependent Legacy Systems:} Early NGO computerization efforts often relied on bespoke local programs managed by a single gifted technician or developer. While functional in the short term, this structure creates existential continuity risks when that individual departs or when transaction volumes outgrow the architecture. Transitioning to documented, multi-user platforms (such as G~Banker) is essential for institutional survival.
\item \textbf{The Last-Mile Verification Gap:} Discrepancies between physical field registers (school attendance books, passbooks, clinic logs) and digital summary reports are pervasive in rural development. When field staff face pressure to demonstrate 100\% achievement, unverified self-reporting easily leads to inflated performance figures unless independent monitoring officers conduct regular register-to-system reconciliation audits.
\item \textbf{Structured Workplans versus Ad-Hoc Reporting:} Effective monitoring cannot rely on sporadic requests from headquarters. Organizations that succeed in maintaining high data quality do so by establishing predictable, pre-scheduled monthly monitoring routines where specific officers are assigned to inspect designated branches and project sites on fixed calendar dates.
\end{enumerate}

% ================= CHAPTER 4 =================
\chapter{Analysis and Findings}\label{chap:analysis}

\chapterglance{Empirical Tools, Pipelines \& Field Gaps}{Analyzes the MEL tool stack, C-RANGPUR digital survey, HCVI vulnerability formulas, G~Banker succession, A2J court docket audits, and last-mile register discrepancies.}

This chapter presents the empirical core of this study. It draws directly on my field observations, system walkthroughs, document reviews, and the running daily log maintained throughout my three-month placement at the RDRS MEL Department.

\section{MIS Tools and Systems Used in Project Monitoring}

\subsection{The observed tool stack}
Inside the MEL Department at Rangpur, data processing is divided across distinct tools based on the nature of the task. Table~\ref{tab:stack} summarizes the primary software applications and manual formats observed.

\begin{table}[htbp]
\centering\small\caption{Operational Tool Stack of the MEL Department}\label{tab:stack}
\begin{tabular}{|>{\raggedright\arraybackslash}p{3.2cm}|>{\raggedright\arraybackslash}p{3.8cm}|>{\raggedright\arraybackslash}p{7.4cm}|}
\hline
\rowcolor{RDRSBlue}\headercell{Monitoring Function} & \headercell{Primary Tool / Platform} & \headercell{Operational Fit \& Justification} \\ \hline
Routine Field Capture & Structured Excel forms \& registers & Standardized templates; auditable paper registers for branch visits \\ \hline
\rowcolor{RDRSLight} Ad-hoc \& Baseline Surveys & KoBo Toolbox (mobile / Enketo) & Dynamic skip logic; offline mobile collection; fast deployment; instant Excel export \\ \hline
Central Compilation & Microsoft Excel & Pivot tables, cross-branch consolidation, and data cleaning \\ \hline
\rowcolor{RDRSLight} Analytical Dashboards & Microsoft Power~BI & Interactive visual reporting, KPI tracking, and donor presentations \\ \hline
Microfinance Operations & G~Banker (South Tech / local) & Centralized web-based MFI platform; single-click POMIS reports; automated PAR \\ \hline
\end{tabular}
\end{table}

For regular branch monitoring, officers rely on structured Excel workbooks aligned with printed field registers. When launching new baseline studies or thematic surveys, the department uses KoBo Toolbox hosted on European servers. Enumerators deploy the KoBo Collect mobile app or web-based Enketo forms on Android devices, allowing offline data collection in flood-prone char areas with automatic synchronization once connectivity is restored.

\subsection{Baseline design lifecycle: lessons from ICPS}\label{sec:icps}
To understand how the MEL Department conceptualizes project measurement from scratch, I studied the baseline design of the ICPS project (\textit{Strengthening Disability-Inclusive Child Protection Systems to Prevent and Respond to Violence against Children with Disabilities in Selected Communities of Bangladesh}). Funded by UNICEF with a budget of BDT~2,20,40,834, ICPS operates under Output 4.2 (social service systems) and Output 4.3 (harmful practices).

During a technical briefing, Md.~Saiful Alam (Shekul) explained that a credible baseline follows a rigorous six-step progression (Figure~\ref{fig:baseline}).

\begin{figure}[htbp]
\centering
\fcolorbox{RDRSTeal}{RDRSLight}{%
\begin{minipage}{0.92\textwidth}
\centering
\scriptsize \fcolorbox{RDRSBlue}{white}{\textbf{\color{RDRSBlue}1. Doc Review}} \textcolor{RDRSBlue}{$\rightarrow$} \fcolorbox{RDRSBlue}{white}{\textbf{\color{RDRSBlue}2. Indicators}} \textcolor{RDRSBlue}{$\rightarrow$} \fcolorbox{RDRSBlue}{white}{\textbf{\color{RDRSBlue}3. MEL Matrix}} \textcolor{RDRSBlue}{$\rightarrow$} \fcolorbox{RDRSBlue}{white}{\textbf{\color{RDRSBlue}4. Concept Paper}} \textcolor{RDRSBlue}{$\rightarrow$} \fcolorbox{RDRSBlue}{white}{\textbf{\color{RDRSBlue}5. Formal TOR}} \textcolor{RDRSBlue}{$\rightarrow$} \fcolorbox{RDRSBlue}{white}{\textbf{\color{RDRSBlue}6. Inception Report}} \\[4pt]
{\scriptsize\color{RDRSGrey} Qualitative Track: FGDs \& KIIs \quad $\bullet$ \quad Quantitative Track: Household Sample Surveys\par}
\end{minipage}}
\caption{The six-step baseline design lifecycle (ICPS project baseline, September 2026).}
\label{fig:baseline}
\end{figure}

UNICEF established 21 specific indicators for ICPS. Qualitative indicators (such as institutional responsiveness and community stigma) are captured via Focus Group Discussions (FGD) and Key Informant Interviews (KII). Quantitative indicators (such as incidence of violence and service uptake) are gathered through household surveys. The Inception Report locks the methodology and field sampling framework before field enumerators begin interviewing.

\subsection{C-RANGPUR baseline survey instrument}\label{sec:crangpur-instrument}
To evaluate digital survey architecture, I conducted a deep dive into the C-RANGPUR baseline survey instrument. The primary reference document was the departmental file \textit{Baseline\_Questionnaires\_of\_CRANGPUR\_-\_Indicator\_Result\_Measerment.docx}, deployed live via KoBo Toolbox (Enketo Express) at \url{https://ee-eu.kobotoolbox.org/x/wzjSgzZ6}.

The digital questionnaire is organized into ten distinct sections (Table~\ref{tab:crangpur-sections}). Sections 1--4 establish household demographic and asset controls; Sections 5--9 directly measure project indicators 1.1 through 3.2; and Section 10 contains the scored Household Climate Vulnerability Index (HCVI).

{\small
\begin{longtable}{|>{\raggedright\arraybackslash}p{0.7cm}|>{\raggedright\arraybackslash}p{3.4cm}|>{\raggedright\arraybackslash}p{6.4cm}|>{\raggedright\arraybackslash}p{2.7cm}|}
\caption{Structure of the C-RANGPUR Baseline Digital Survey Instrument}\label{tab:crangpur-sections}\\
\hline
\rowcolor{RDRSBlue}\headercell{\#} & \headercell{Section Title} & \headercell{Core Content \& Variables Captured} & \headercell{Target Indicator} \\ \hline
\endfirsthead
\hline
\rowcolor{RDRSBlue}\headercell{\#} & \headercell{Section Title} & \headercell{Core Content \& Variables Captured} & \headercell{Target Indicator} \\ \hline
\endhead
\hline
\endfoot
\hline
\endlastfoot
1 & Location Background & District, Upazila, Union, Ward, Village; Hazard zones: Char, Riverbank erosion, Flood, Waterlogging, Drought & Geographic controls \\ \hline
\rowcolor{RDRSLight} 2 & Respondent Profile & Name, sex, NID, phone, ethnicity, marital status, education, occupation, disability type, prior RDRS membership & Beneficiary identity \\ \hline
3 & Household Composition & Household head demographics; counts of children under 5, pregnant/lactating women, elderly (60+), disabled & Vulnerability profile \\ \hline
\rowcolor{RDRSLight} 4 & Assets \& Infrastructure & Housing material, electricity, drinking water, latrines, 14 livestock/asset categories, land tenure & Baseline asset base \\ \hline
5 & Livelihoods \& Income & Income sources (15 codes), IGA support, CSA awareness/adoption (18 practices), 3-year income trend & Indicators 1.1, 1.2 \\ \hline
\rowcolor{RDRSLight} 6 & Community Participation & Membership in 21 community institutions, meeting attendance, voice in decision-making, female leadership & Indicator 2.1 \\ \hline
7 & Social Services \& Claims & Safety-net receipt (14 types), claims filed, official response received, resolution status & Indicator 2.2 \\ \hline
\rowcolor{RDRSLight} 8 & Primary Health Access & Health-seeking behaviour, choice across 13 service facilities (incl.\ RDRS Eye Care), barriers to treatment & Indicator 3.1 \\ \hline
9 & Non-Communicable Diseases & 8 diagnosed NCD conditions (hypertension, diabetes, cataract, etc.), screening history, lifestyle habits & Indicator 3.2 \\ \hline
\rowcolor{RDRSLight} 10 & Climate Vulnerability & 16-item scored module across Exposure (5), Sensitivity (5), and Adaptive Capacity (6) domains & HCVI composite \\ \hline
\end{longtable}
}

\subsection{Indicator formulas and mathematical definitions}\label{sec:crangpur-indicators}
In the project monitoring matrix (file \textit{F01-112}), each indicator is mathematically defined with explicit numerators and denominators.

\formulacard[RDRSBlue]{Standard Indicator Calculation Logic}{%
$\displaystyle \text{Indicator Result (\%)} = \left(\frac{\text{Observed Target Numerator (Eligible Population Meeting Criteria)}}{\text{Total Baseline Denominator (Surveyed Sample Group)}}\right) \times 100$
}{%
Standard percentage formulation applied across all C-RANGPUR indicators in the MEL matrix F01-112.
}

Table~\ref{tab:crangpur-indicators} details the exact mathematical formulas, operational definitions, and target thresholds established for C-RANGPUR.

{\small
\begin{longtable}{|>{\centering\arraybackslash}p{1.8cm}|>{\raggedright\arraybackslash}p{4.4cm}|>{\raggedright\arraybackslash}p{4.8cm}|>{\centering\arraybackslash}p{2.4cm}|}
\caption{Mathematical Definitions of C-RANGPUR Project Indicators}\label{tab:crangpur-indicators}\\
\hline
\rowcolor{RDRSBlue}\headercell{Indicator} & \headercell{Operational Definition} & \headercell{Numerator / Denominator Specification} & \headercell{Formula} \\ \hline
\endfirsthead
\hline
\rowcolor{RDRSBlue}\headercell{Indicator} & \headercell{Operational Definition} & \headercell{Numerator / Denominator Specification} & \headercell{Formula} \\ \hline
\endhead
\hline
\endfoot
\hline
\endlastfoot
1.1 & Climate-affected households adopting $\ge$1 additional IGA (+25\% from baseline) & Num: Households adopting $\ge$1 new IGA \newline Den: Total targeted households surveyed & $\frac{\text{Num}}{\text{Den}} \times 100$ \\ \hline
\rowcolor{RDRSLight} 1.2 & $\ge$40\% of farming households increasing income by 20\% via Climate Smart Agriculture & Num: Households adopting $\ge$1 promoted CSA practice \newline Den: Total targeted farming households surveyed & $\frac{\text{Num}}{\text{Den}} \times 100$ \\ \hline
2.1 & Active participation of disadvantaged people in local decision-making (+30\% from baseline) & Num: Disadvantaged members attending meetings, speaking, AND having opinions considered \newline Den: Total disadvantaged respondents surveyed & $\frac{\text{Num}}{\text{Den}} \times 100$ \\ \hline
\rowcolor{RDRSLight} 2.1 (women) & Female active participation in community decision-making bodies & Num: Female respondents meeting all 2.1 criteria \newline Den: Total female respondents surveyed & $\frac{\text{Num}}{\text{Den}} \times 100$ \\ \hline
2.2 & Proportion of citizen claims successfully resolved by local public authorities & Num: Claims fully or partially resolved by authorities \newline Den: Total formal claims submitted in period & $\frac{\text{Num}}{\text{Den}} \times 100$ \\ \hline
\rowcolor{RDRSLight} 3.1 & 25\% of disadvantaged members accessing formal primary health care & Num: Respondents receiving care at public/RDRS clinic \newline Den: Respondents reporting health need in last 12 months & $\frac{\text{Num}}{\text{Den}} \times 100$ \\ \hline
3.2 & Reduction in selected NCD prevalence by 25\% against baseline & Num: Adults diagnosed with $\ge$1 selected NCD \newline Den: Total adult respondents surveyed & $\frac{\text{Num}}{\text{Den}} \times 100$ \\ \hline
\end{longtable}
}

In Section 7, claims tracking is further decomposed into five sub-metrics: total claims submitted, total claims acknowledged by authorities, total claims resolved, response rate ($[\text{Responded} \div \text{Submitted}] \times 100$), and resolution rate ($[\text{Resolved} \div \text{Submitted}] \times 100$).

\subsection{Household Climate Vulnerability Index (HCVI)}\label{sec:hcvi}
The most sophisticated analytical tool developed by the MEL Department is the Household Climate Vulnerability Index (HCVI), captured in Section 10 of the questionnaire. The index builds on the international climate vulnerability framework (IPCC, 2022; Hahn et al., 2009; FAO, 2016):
\begin{equation}
\mathrm{Vulnerability} = \mathrm{Exposure} + \mathrm{Sensitivity} - \mathrm{Adaptive\ Capacity}
\label{eq:hcvi-logic}
\end{equation}

\formulacard[RDRSTeal]{Household Climate Vulnerability Index (HCVI) Master Composite Formula}{%
$\displaystyle \mathbf{HCVI = \underbrace{(\mathrm{EI} \times 0.30)}_{\text{30\% Exposure Weight}} + \underbrace{(\mathrm{SI} \times 0.30)}_{\text{30\% Sensitivity Weight}} + \underbrace{\bigl[(100 - \mathrm{ACI}) \times 0.40\bigr]}_{\text{40\% Adaptive Capacity Deficit Weight}}}$
}{%
Composite index bounded between 0 and 100. Lower values reflect climate resilience; higher values indicate severe vulnerability.
}

To calculate the composite score, the three sub-indices are normalized to a standard 0--100 percentage scale:

\vspace{0.15cm}
\begin{center}
\setlength{\tabcolsep}{2pt}
\begin{tabular}{ccc}
\fcolorbox{RDRSBlue}{white}{\parbox{4.3cm}{\centering \scriptsize\textbf{\color{RDRSBlue}EXPOSURE INDEX (EI)}\\[3pt]\small $\displaystyle \mathrm{EI} = \left(\frac{\text{Raw Score}}{20}\right) \times 100$\\[3pt]\tiny 5 items $\times$ 4 pts max = 20 pts}} &
\fcolorbox{RDRSBlue}{white}{\parbox{4.3cm}{\centering \scriptsize\textbf{\color{RDRSBlue}SENSITIVITY INDEX (SI)}\\[3pt]\small $\displaystyle \mathrm{SI} = \left(\frac{\text{Raw Score}}{20}\right) \times 100$\\[3pt]\tiny 5 items $\times$ 4 pts max = 20 pts}} &
\fcolorbox{RDRSBlue}{white}{\parbox{4.3cm}{\centering \scriptsize\textbf{\color{RDRSBlue}ADAPTIVE CAPACITY (ACI)}\\[3pt]\small $\displaystyle \mathrm{ACI} = \left(\frac{\text{Raw Score}}{24}\right) \times 100$\\[3pt]\tiny 6 items $\times$ 4 pts max = 24 pts}}
\end{tabular}
\end{center}
\vspace{0.15cm}

The scoring components within each domain are structured as follows:
\begin{itemize}
\item \textbf{Domain A: Exposure (Q10.1--10.5, Max 20 points):} Frequency of climate hazards in last 3 years (0--4); number of times affected (0--4); months per year affected (0--4); temporary evacuations (0--4); daily livelihood disruptions (0--4).
\item \textbf{Domain B: Sensitivity (Q10.6--10.10, Max 20 points):} Months of food shortage in last year (0--4); duration of income loss during last disaster (0--4); physical asset damage level (0--4); safe drinking water crisis severity (0--4); number of dependent vulnerable members needing emergency care (0--4).
\item \textbf{Domain C: Adaptive Capacity (Q10.11--10.16, Max 24 points):} Early warning reception (0--4); disaster preparedness measures taken (0--4); adoption of $\ge$1 CSA practice (No=0 / Yes=4); adoption of $\ge$1 new risk-reducing IGA (No=0 / Yes=4); community group participation (0--4); preventive health measures taken (No=0 / Yes=4).
\end{itemize}

\begin{table}[htbp]
\centering\small\caption{HCVI Score Interpretation Bands}\label{tab:hcvi-bands}
\begin{tabular}{|>{\raggedright\arraybackslash}p{2.0cm}|>{\raggedright\arraybackslash}p{2.6cm}|>{\raggedright\arraybackslash}p{2.6cm}|>{\raggedright\arraybackslash}p{2.8cm}|>{\raggedright\arraybackslash}p{3.6cm}|}
\hline
\rowcolor{RDRSBlue}\headercell{Score Band} & \headercell{EI Level} & \headercell{SI Level} & \headercell{ACI Level} & \headercell{HCVI Vulnerability Status} \\ \hline
0--$<$20 & Very Low Exposure & Very Low Sensitivity & Very Low Capacity & \textcolor{RDRSTeal}{\textbf{Very Low Vulnerability}} \\ \hline
\rowcolor{RDRSLight} 20--$<$40 & Low Exposure & Low Sensitivity & Low Capacity & \textcolor{RDRSTeal}{\textbf{Low Vulnerability}} \\ \hline
40--$<$60 & Moderate Exposure & Moderate Sensitivity & Moderate Capacity & \textcolor{RDRSBlue}{\textbf{Moderate Vulnerability}} \\ \hline
\rowcolor{RDRSLight} 60--$<$80 & High Exposure & High Sensitivity & High Capacity & \textcolor{RDRSOrange}{\textbf{High Vulnerability}} \\ \hline
80--100 & Very High Exposure & Very High Sensitivity & Very High Capacity & \textcolor{red}{\textbf{Very High Vulnerability}} \\ \hline
\end{tabular}
\end{table}

As Md.~Ajazul Haque explained during our walkthrough, Exposure and Sensitivity represent external environmental hazards that RDRS cannot directly stop (a family living on a char cannot easily change the river's path). Adaptive Capacity represents community resilience---the exact dimension where RDRS training, credit, and agricultural inputs intervene.

\subsection{The MEL matrix and the C-RANGPUR digital pipeline}\label{sec:mel-matrix}
The foundational planning tool for project monitoring is the MEL Matrix, maintained in spreadsheet file \textit{F01-112} under the tab ``Monitoring Planning''. Table~\ref{tab:mel-11} shows the worked planning matrix for Indicator 1.1.

\begin{table}[htbp]
\centering\small\caption{MEL Matrix Worked Planning Specification: Indicator 1.1}\label{tab:mel-11}
\begin{tabular}{|>{\raggedright\arraybackslash}p{3.4cm}|>{\raggedright\arraybackslash}p{9.8cm}|}
\hline
\rowcolor{RDRSBlue}\headercell{Matrix Field} & \headercell{Operational Specification in File F01-112} \\ \hline
Indicator Title & Targeted climate-affected households adopting $\ge$1 additional IGA, +25\% from baseline \\ \hline
\rowcolor{RDRSLight} Quantitative Data Needed & Household ID; baseline/current IGAs; enterprise type; launch date; monthly income earned \\ \hline
Qualitative Data Needed & Motivation for adoption; operational challenges; market barriers; perceived household benefits \\ \hline
\rowcolor{RDRSLight} Disaggregation Dimensions & Sex and age of household head; disability status; geographic zone; poverty classification \\ \hline
Data Source & Target rural households in program unions \\ \hline
\rowcolor{RDRSLight} Collection Method & Structured household survey; individual Key Informant Interviews \\ \hline
Reporting Schedule & Baseline survey, semi-annual progress review, and annual consolidation \\ \hline
\end{tabular}
\end{table}

Figure~\ref{fig:crangpur-flow} illustrates the end-to-end data pipeline for C-RANGPUR as detailed in my technical interview with Md.~Ajazul Haque.

\begin{figure}[htbp]
\centering
\fcolorbox{RDRSBlue}{RDRSLight}{%
\begin{minipage}{0.92\textwidth}
\centering
\small\fcolorbox{RDRSBlue}{white}{\textbf{\color{RDRSBlue}KoBo Mobile Capture}} \textcolor{RDRSBlue}{$\rightarrow$} \small\fcolorbox{RDRSTeal}{white}{Cloud Server} \textcolor{RDRSBlue}{$\rightarrow$} \small\fcolorbox{RDRSBlue}{white}{\textbf{\color{RDRSBlue}MEL Data Verification}} \textcolor{RDRSBlue}{$\rightarrow$} \small\fcolorbox{RDRSTeal}{white}{Excel Cleaning} \textcolor{RDRSBlue}{$\rightarrow$} \small\fcolorbox{RDRSOrange}{white}{\textbf{\color{RDRSOrange}Power~BI Dashboard}} \\[3pt]
{\scriptsize\color{RDRSGrey} HCVI and IGA indicators calculated across tools \quad $\bullet$ \quad Completely isolated from G~Banker MFI database\par}
\end{minipage}}
\caption{The C-RANGPUR data pipeline: from mobile survey capture to Power~BI visualization.}
\label{fig:crangpur-flow}
\end{figure}

Table~\ref{tab:kii-ajazul} summarizes the key technical clarifications obtained during this interview.

\begin{table}[htbp]
\centering\small\caption{Key Informant Interview: Md.~Ajazul Haque, Manager (MEL)}\label{tab:kii-ajazul}
\begin{tabular}{|>{\centering\arraybackslash}p{0.8cm}|>{\raggedright\arraybackslash}p{4.2cm}|>{\raggedright\arraybackslash}p{8.2cm}|}
\hline
\rowcolor{RDRSBlue}\headercell{\#} & \headercell{Technical Inquiry} & \headercell{Manager's Response \& Operational Finding} \\ \hline
1 & Is the HCVI composite calculation automated inside KoBo? & No single automated tool exists; raw survey data exports to Excel where weighting macros run before loading to Power~BI. \\ \hline
\rowcolor{RDRSLight} 2 & What is the function of the numeric ``Code'' column in the survey? & Codes support export into SPSS for advanced statistical correlation and variance analysis. \\ \hline
3 & Does KoBo submission update dashboards in real time? & Data auto-submits to the KoBo server, but passes a manual MEL verification and cleaning gate before entering Power~BI. \\ \hline
\rowcolor{RDRSLight} 4 & How are baseline values tracked against end-of-project targets? & Baseline values establish the locked minimum threshold; routine monitoring surveys compare annual progress against that fixed benchmark. \\ \hline
5 & Is project MEL integrated with microfinance software (G~Banker)? & The two systems are completely separated; development M\&E runs on KoBo/Power~BI while microfinance runs exclusively on G~Banker. \\ \hline
\end{tabular}
\end{table}

\subsection{Microfinance software evolution: Micro Bank to South Tech to G~Banker}
The evolution of microfinance software at RDRS offers a valuable case study in system continuity. During our discussions, Mr.~Iftekharul Islam (Manager-Monitoring) traced this thirty-year journey.

In 1995--1997, when RDRS formalized its microfinance operations, field officers were called ATMs (Assistant Thana Managers). The organization developed its first computerized management system, known as Micro Bank, supported by technical assistance from FAO and GTZ. A dedicated team of 27 programmers built the platform, including my grandfather (Boro Abbu) alongside a senior programmer named Mahabub. Micro Bank successfully digitized branch ledgers, but as transaction volumes expanded, reporting became overly dependent on a handful of individuals who understood the underlying database queries. If a key programmer was unavailable, month-end consolidation stalled.

\begin{figure}[htbp]
\centering
\fcolorbox{RDRSBlue}{RDRSLight}{%
\begin{minipage}{0.92\textwidth}
\centering
{\scriptsize\color{RDRSOrange}\textbf{ERA 1: PERSON-DEPENDENT}} \hspace{0.4cm} {\scriptsize\color{RDRSTeal}\textbf{ERA 2: TRANSITIONAL}} \hspace{0.4cm} {\scriptsize\color{RDRSBlue}\textbf{ERA 3: ENTERPRISE AUTOMATION}} \\[3pt]
\small\fcolorbox{RDRSOrange}{white}{\parbox{3.6cm}{\centering \textbf{Micro Bank}\\ \scriptsize 1995--1997, FAO/GTZ\\ \scriptsize 27 programmers\\ \scriptsize High key-person risk}} \textcolor{RDRSBlue}{$\longrightarrow$} \small\fcolorbox{RDRSTeal}{white}{\parbox{3.2cm}{\centering \textbf{South Tech}\\ \scriptsize 2011 rollout\\ \scriptsize Commercial database\\ \scriptsize Reduced person risk}} \textcolor{RDRSBlue}{$\longrightarrow$} \small\fcolorbox{RDRSBlue}{white}{\parbox{3.6cm}{\centering \textbf{\color{RDRSBlue}G~Banker}\\ \scriptsize 2017/18--Present\\ \scriptsize Centralized web system\\ \scriptsize 1-click POMIS reporting}} \\
\end{minipage}}
\caption{The three eras of microfinance software at RDRS Bangladesh: 1995 to 2026.}
\label{fig:missuccession}
\end{figure}

To eliminate this operational bottleneck, RDRS partnered with South Tech in 2011 to implement a commercial microfinance package. My grandfather assisted with data migration and system configuration before retiring. Between 2017 and 2018, RDRS upgraded to G~Banker, a modern web-based microfinance enterprise solution developed in Bangladesh.

Today, G~Banker operates across all 288 branches. Routine statements that once took days of manual compilation now generate instantaneously:
\begin{itemize}
\item \textbf{CO-wise Program Changes Report:} Displays monthly changes per Credit Organizer (active samities, borrower counts, disbursements, recoveries, overdue balances, write-offs, and PAR).
\item \textbf{POMIS-5 (Statement of Financial Position):} Branch-level balance sheet comparing current balances against the prior financial year.
\item \textbf{POMIS-4 (Statement of Comprehensive Income and Expenditure):} Monthly and annual branch profit-and-loss statements generated directly from daily ledger entries.
\end{itemize}

\section{Data Flow: From Grassroots Field Units to Management Dashboards}
Field data reaches central management through two distinct reporting channels, coordinated via the monthly departmental monitoring schedule (Figure~\ref{fig:dataflow}).

\begin{figure}[htbp]
\centering
\fcolorbox{RDRSBlue}{RDRSLight}{%
\begin{minipage}{0.92\textwidth}
\centering
\small\fcolorbox{RDRSBlue}{white}{\textbf{\color{RDRSBlue}Field Visits}} \textcolor{RDRSBlue}{$\rightarrow$} \small\fcolorbox{RDRSTeal}{white}{Paper Registers \& KoBo} \textcolor{RDRSBlue}{$\rightarrow$} \small\fcolorbox{RDRSBlue}{white}{\textbf{\color{RDRSBlue}MEL Branch Audits}} \textcolor{RDRSBlue}{$\rightarrow$} \small\fcolorbox{RDRSOrange}{white}{\textbf{\color{RDRSOrange}Executive MIS Dashboards}} \\[3pt]
{\scriptsize\color{RDRSGrey} Routine microfinance: bi-monthly branch visits \quad $\bullet$ \quad Development projects: monthly indicator schedules\par}
\end{minipage}}
\caption{The dual-track field monitoring flow coordinated by the MEL Department.}
\label{fig:dataflow}
\end{figure}

The September 2026 monitoring schedule illustrates how departmental duties are structured:
\begin{itemize}
\item \textbf{Microfinance Track:} Ten recurring supervisory tasks assigned to designated officers visiting branches twice monthly (02--15 Sept and 16--28 Sept). Tasks include cross-checking borrower passbooks against samity registers, verifying cash collections, auditing loan disbursement vouchers, inspecting overdue recovery plans, and verifying write-off documentation.
\item \textbf{Development Project Track:} Time-bound indicator tracking across active donor grants. In September 2026, this covered C-RANGPUR baseline data collection across six districts; MCGL health data validation in Panchagarh; Jononni maternal health DQA in Lalmonirhat; SEEDS education monitoring in Rangpur; and the live ICPS baseline across Mymensingh, Rajshahi, and Rangpur Divisions.
\end{itemize}

\section{Special Monitoring Protocols: Where MIS Meets Physical Verification}\label{sec:special-monitoring}
A core insight from my placement is that digital dashboards cannot replace physical inspection. The MEL Department enforces a strict ``special monitoring'' checklist:
\begin{enumerate}
\item \textbf{Overdue Loan Investigations:} When G~Banker flags a delinquent account, an officer must physically visit the borrower's home, record a written reason for default, and submit that finding to the branch loan committee.
\item \textbf{Disbursement Paper Audits:} Comparing loan amounts approved by committees against actual cash disburse\-ment receipts and passbook entries to detect phantom loans.
\item \textbf{High-Value Tracking:} Independent monthly tracking of high-exposure accounts using Portfolio-at-Risk (PAR) and Overdue-to-Total-Receivables (OTR) metrics.
\item \textbf{Misappropriation and Fraud Tracking:} Dedicated investigation procedures for suspected cash skimming, with formal grievance boxes installed at every branch office.
\end{enumerate}

Software identifies anomalies, but human investigators verify the ground reality. The strength of RDRS's system lies in this division of labor.

\section{Risk Mitigation Systems: Credit Shield Microinsurance}
Since July 2024, RDRS has operated Credit Shield Insurance in partnership with Guardian Life Insurance Limited. Microcredit borrowers pay a mandatory premium of BDT~6 per BDT~1,000 borrowed. In the event of a borrower's or co-borrower's death, the policy waives the entire outstanding loan balance (principal plus accrued service charges up to BDT~20 lakh) and provides an immediate cash funeral grant of BDT~15,000 to the bereaved family.

From an MIS perspective, the claim workflow reveals a paper-heavy bottleneck:
\begin{center}
\small\textbf{Death Confirmed at Branch} $\longrightarrow$ \textbf{Branch Focal Verifies Dossier (2 days)} $\longrightarrow$ \textbf{Regional Focal Review} $\longrightarrow$ \textbf{Central Focal (Dhaka)} $\longrightarrow$ \textbf{Guardian Life Review (10 days)} $\longrightarrow$ \textbf{Disbursement (5 days)}
\end{center}
While G~Banker records the eventual loan waiver, all supporting documents (death certificates, medical records, union parishad attestations) travel physically as paper dossiers between field branches and Dhaka. Digitizing this claim pipeline represents a significant efficiency opportunity.

\section{Regulatory Compliance Data Systems: AML/CFT}
Under the Anti-Money Laundering (AML) and Combating the Financing of Terrorism (CFT) Policy (effective January 2022), RDRS maintains an internal compliance database tracking senior leaders, trustees, and high-value transactions. This system captures National ID numbers, tax identifiers (TIN), permanent addresses, and affiliated board positions, generating statutory filings for Bangladesh Bank and the Bangladesh Financial Intelligence Unit (BFIU). All cash transactions exceeding BDT~5,000,000 require automated compliance flagging.

\section{Credit Governance and Approval Authority}
Loan approval limits follow a clear hierarchical ladder:
\begin{center}
\small\textbf{Credit Organizer (Recommends)} $\longrightarrow$ \textbf{Branch Manager ($\le$BDT 1.5L)} $\longrightarrow$ \textbf{Area Manager ($\le$BDT 3.0L)} $\longrightarrow$ \textbf{Regional Manager ($\le$BDT 5.0L)} $\longrightarrow$ \textbf{Head of Microfinance ($>$BDT 5.0L)}
\end{center}
Monitoring officers trace this hierarchy in reverse during branch audits: confirming that large loans carry appropriate managerial signatures before disbursement.

\section{Case Study: Access to Justice for Women (A2J)}\label{sec:a2j}
To evaluate how monitoring functions in legal and social sectors, I reviewed the Access to Justice for Women (A2J) project \cite{a2jdoc}, implemented jointly with GIZ and the Law and Justice Division (Ministry of Law, Justice and Parliamentary Affairs).

\subsection{Institutional split: formal versus informal justice}
A2J operates through a unique institutional partnership in Rangpur District:
\begin{itemize}
\item \textbf{Formal Justice System (RDRS Bangladesh):} RDRS paralegals operate inside the formal court complex, District Legal Aid Office (DLAO), police stations, and Rangpur Central Prison, managing formal litigation and court-annexed mediation.
\item \textbf{Informal Justice System (Nijera Kori):} Partners facilitate Community Dispute Resolution (CDR) and village-level mediation.
\end{itemize}

Through April 2025, the project assisted 6,092 litigants (3,348 female, 2,744 male) and established a permanent Legal Service Information Desk (LSID) inside the Rangpur District Court. On 18 March 2025, the Senior District and Sessions Judge approved an official panel of trained mediators, institutionalizing court-annexed alternative dispute resolution (ADR) for compoundable offenses (\textit{aposh-joggo}) under the Code of Criminal Procedure.

\subsection{The April 2025 physical court docket audit}
Because Bangladesh's subordinate courts lack an integrated digital query system accessible to civil society, RDRS paralegals conducted a comprehensive physical docket audit in April 2025. Paralegals manually inspected courtroom registers (\textit{Ejlash}) and administrative clerk archives (\textit{Serestha}) across every bench in Rangpur.

\begin{figure}[htbp]
\centering
\fcolorbox{RDRSBlue}{RDRSLight}{%
\begin{minipage}{0.92\textwidth}
\centering\textbf{\color{RDRSBlue}Rangpur District Court Caseload at a Glance (April 2025 Physical Audit)}\\[4pt]
\begin{tabular}{p{3.8cm}p{7.8cm}}
\textbf{Total Active Cases:} & \textbf{77,849 cases counted across all benches} \\[2pt]
\textbf{Criminal Cases (56.8\%):} & \textbf{44,215 files} {\color{RDRSBlue}\rule{4.8cm}{3mm}} \\
$\bullet$ Judicial Magistracy: & 15,794 files {\color{RDRSTeal}\rule{3.4cm}{2.4mm}} \\
$\bullet$ Metro Magistracy: & 9,818 files {\color{RDRSTeal}\rule{2.1cm}{2.4mm}} \\
$\bullet$ District Sessions: & 8,936 files {\color{RDRSTeal}\rule{1.9cm}{2.4mm}} \\
$\bullet$ Nari O Shishu Tribunals: & 6,068 files {\color{RDRSOrange}\rule{1.3cm}{2.4mm}} (heavy domestic violence docket) \\
$\bullet$ Metropolitan Sessions: & 2,403 files {\color{RDRSTeal}\rule{0.5cm}{2.4mm}} \\
$\bullet$ Other Tribunals: & 1,196 files {\color{RDRSTeal}\rule{0.25cm}{2.4mm}} (ATT, CT, AHTT, ST) \\[4pt]
\textbf{Civil Cases (43.2\%):} & \textbf{33,634 files} {\color{RDRSBlue}\rule{3.6cm}{3mm}} \\
$\bullet$ Land Survey Tribunal: & 5,493 files {\color{RDRSOrange}\rule{1.2cm}{2.4mm}} (highest single civil load) \\
\end{tabular}\\[4pt]
{\scriptsize\color{RDRSGrey} Data manually enumerated by 6 A2J paralegals inspecting \textit{Ejlash} (bench registers) and \textit{Serestha} (clerk archives).\par}
\end{minipage}}
\caption{Distribution of active court cases in Rangpur District (April 2025 physical audit).}
\label{fig:court-caseload}
\end{figure}

The physical audit documented 77,849 pending cases across Rangpur District: 44,215 criminal cases (56.8\%) and 33,634 civil disputes (43.2\%). Table~\ref{tab:court-audit} breaks down the criminal caseload by jurisdiction.

\begin{table}[htbp]
\centering\small\caption{Criminal Caseload Breakdown: Rangpur District Courts (April 2025 Audit)}\label{tab:court-audit}
\begin{tabular}{|c|>{\raggedright\arraybackslash}p{2.4cm}|>{\raggedright\arraybackslash}p{6.4cm}|r|r|}
\hline
\rowcolor{RDRSBlue}\headercell{Sl} & \headercell{Jurisdiction} & \headercell{Court / Tribunal Category} & \headercell{Sub-total} & \headercell{Total Files} \\ \hline
1 & Magistracy & Judicial Magistracy (including GRO dockets) & 15,794 & 25,612 \\
  & & Metropolitan Magistracy (including GRO dockets) & 9,818 & \\ \hline
\rowcolor{RDRSLight} 2 & Sessions Courts & District Sessions Bench & 8,936 & 11,339 \\
\rowcolor{RDRSLight}   & & Metropolitan Sessions Bench & 2,403 & \\ \hline
3 & Tribunals & Nari O Shishu Nirjatan Daman Tribunals & 6,068 & 7,264 \\
  & & Specialized Tribunals (ATT, CT, AHTT, ST) & 1,196 & \\ \hline
\multicolumn{3}{|l|}{\textbf{Total Criminal Cases Counted}} & \multicolumn{2}{r|}{\textbf{44,215}} \\ \hline
\end{tabular}
\end{table}

This audit underscores the immense manual reporting burden placed on field staff: six paralegals had to manually count tens of thousands of paper dockets while simultaneously providing legal advice to indigent clients.

\section{Observed Strengths of Current MIS Practices}
RDRS demonstrates several clear strengths in its monitoring setup:
\begin{itemize}
\item \textbf{Single-Click MFI Reporting:} G~Banker delivers immediate branch financial statements and CO performance metrics, saving hundreds of hours of monthly manual compilation.
\item \textbf{Disciplined Monthly Schedules:} The MEL Department operates from a transparent monthly calendar that explicitly links named officers to specific branches and projects.
\item \textbf{Appropriate Tool Selection:} The division between Excel (for stable, standardized routines) and KoBo Toolbox (for rapid, offline mobile survey deployment) fits field realities.
\end{itemize}

\section{Observed Gaps and Field-Level Data Discrepancies}
Despite strong digital platforms, significant vulnerabilities persist at the point of field capture.

\begin{figure}[htbp]
\centering
\fcolorbox{RDRSBlue}{RDRSLight}{%
\begin{minipage}{0.92\textwidth}
\centering\textbf{\color{RDRSBlue}Last-Mile Data Integrity Gaps (Somriddhi Program Audit, May 2023)}\\[4pt]
\begin{tabular}{p{3.4cm}p{8.2cm}}
Student Attendance: & MIS entered: 100\% {\color{RDRSBlue}\rule{4.0cm}{3mm}} \newline Field register actual: 89\% {\color{RDRSOrange}\rule{3.56cm}{3mm}} \\[3pt]
Rehabilitation Assets: & Valid records found: 35\% {\color{RDRSTeal}\rule{1.4cm}{3mm}} \newline Missing expense vouchers: 65\% {\color{RDRSOrange}\rule{2.6cm}{3mm}} \\[3pt]
Tree Survival Rates: & Debiduba union: 49\% {\color{RDRSTeal}\rule{1.96cm}{3mm}} \newline Beruabari union: 35\% {\color{RDRSOrange}\rule{1.4cm}{3mm}} \\
\end{tabular}\\[4pt]
{\scriptsize\color{RDRSGrey} Archival review of PKSF-funded Somriddhi field monitoring report.\par}
\end{minipage}}
\caption{Visual illustration of field-versus-MIS discrepancies observed in audit archives.}
\label{fig:somriddhi}
\end{figure}

An analysis of the May 2023 field monitoring report on the Somriddhi Programme (PKSF-funded, covering Panchagarh and Kurigram) revealed critical last-mile discrepancies (Figure~\ref{fig:somriddhi}):
\begin{itemize}
\item \textbf{Attendance Inflation:} In Beruabari union, central MIS reports recorded 100\% student attendance in non-formal education centers, but physical inspections of classroom registers showed actual attendance was only 89\%.
\item \textbf{Missing Expenditure Records:} For ultra-poor rehabilitation asset grants (BDT~100,000 per household), 65\% of recipient expenditure receipts could not be traced in field files.
\item \textbf{High Staff Turnover:} 11 out of 74 learning center teachers resigned within a four-month period, creating severe reporting disruptions that were not reflected in central KPI summaries.
\end{itemize}

In C-RANGPUR, I observed minor but telling survey inconsistencies: the project title is spelled ``C RANGPU'' in question 5.3.2 and ``C RANGPUR'' in question 5.4.1; several multi-select questions collapse into single codes in Word exports; and the HCVI vulnerability index requires error-prone manual file transfers between KoBo, Excel, and Power~BI.

\section{Summary Findings Matrix}
Table~\ref{tab:findings} consolidates the operational findings, tool assessments, and data strengths and gaps identified across this evaluation.

\begin{table}[htbp]
\centering\small\caption{Consolidated Evaluation Findings: RDRS MIS and Monitoring Practices}\label{tab:findings}
\begin{tabular}{|>{\raggedright\arraybackslash}p{3.2cm}|>{\raggedright\arraybackslash}p{5.6cm}|>{\raggedright\arraybackslash}p{4.4cm}|}
\hline
\rowcolor{RDRSBlue}\headercell{MIS System / Function} & \headercell{Field Observation \& Current Workflow} & \headercell{Operational Status} \\ \hline
G~Banker Platform & Instant generation of CO-wise Changes, POMIS-4, and POMIS-5 statements & \strengthtag{} High speed, consistency, and automated PAR calculation \\ \hline
\rowcolor{RDRSLight} Monthly Monitoring Plan & Transparent monthly schedule assigning named officers, branches, and dates & \strengthtag{} Clear accountability and structured field coverage \\ \hline
KoBo Mobile Tools & Digital survey deployment with GPS capture and offline mobile capability & \strengthtag{} Fast deployment; eliminates manual paper survey re-entry \\ \hline
\rowcolor{RDRSLight} Register-to-MIS Accuracy & Field audit revealed 89\% register attendance reported as 100\% in central MIS & \gaptag{} Critical last-mile data mismatch; lack of reconciliation gate \\ \hline
Livelihood Asset Records & 65\% of asset grant expense vouchers missing during Somriddhi field inspections & \gaptag{} Weak paper record archiving at peripheral branch level \\ \hline
\rowcolor{RDRSLight} C-RANGPUR HCVI Pipeline & Vulnerability calculation divided across KoBo, manual Excel, and Power~BI & \gaptag{} Manual file handoffs; lack of automated pipeline \\ \hline
Credit Shield Insurance & Paper-based claim files transported through 4 administrative tiers & \gaptag{} Slow physical dossier routing; settlement lag \\ \hline
\rowcolor{RDRSLight} A2J Court Docket Tracking & Six paralegals manually counted 77,849 active cases across court archives & \gaptag{} Heavy manual counting burden diverting staff from legal aid \\ \hline
Microfinance Continuity & Successful transition from legacy Micro Bank to commercial G~Banker & \strengthtag{} Overcame single-person dependency; institutionalized rules \\ \hline
\end{tabular}
\end{table}

\section{Placement Activity Log}\label{sec:dailylog}
To document the empirical foundation of this study, Table~\ref{tab:log} details the chronological log of field visits, staff walkthroughs, system audits, and departmental rotations completed during my placement. This sequence directly reflects the official internship schedule approved by the Head of MEL (Appendix~\ref{sec:official-schedule}).

{\small
\begin{longtable}{|>{\raggedright\arraybackslash}p{1.8cm}|>{\raggedright\arraybackslash}p{3.8cm}|>{\raggedright\arraybackslash}p{3.6cm}|>{\raggedright\arraybackslash}p{5.0cm}|}
\caption{Comprehensive Internship Activity and Fieldwork Log}\label{tab:log}\\
\hline
\rowcolor{RDRSBlue}\headercell{Date / Period} & \headercell{Placement Activity} & \headercell{Staff Contact / Source} & \headercell{Key MIS Insights \& Takeaways} \\ \hline
\endfirsthead
\hline
\rowcolor{RDRSBlue}\headercell{Date / Period} & \headercell{Placement Activity} & \headercell{Staff Contact / Source} & \headercell{Key MIS Insights \& Takeaways} \\ \hline
\endhead
\hline
\endfoot
\hline
\endlastfoot
01 Sep 2026 & Departmental Induction \& Reporting Guidelines & Mr.~Mamunur Rashid (Head of MEL) \newline Dabasish Mahata (Coordinator) & Reviewed RDRS 6th Strategic Plan; mapped institutional reach: 406 Federations, 33 districts, 1.64M participants. \\ \hline
\rowcolor{RDRSLight} 02--04 Sep 2026 & Overview of RDRS Programs \& MIS Toolchains & Mr.~Iftekharul Islam; Fahmida Mostafa (Senior MEL Officer) & Mapped three-tier data tools: physical registers/Excel for routine monitoring; KoBo on mobile; Power~BI dashboards. \\ \hline
07--09 Sep 2026 & Historical Walkthrough of Microfinance Software & Mr.~Iftekharul Islam (Manager-Monitoring) & Traced 3 eras: manual ATMs $\to$ Micro Bank (1995--97, grandfather's team) $\to$ South Tech (2011) $\to$ G~Banker (2017/18). \\ \hline
\rowcolor{RDRSLight} 10--11 Sep 2026 & Audit of Live G~Banker Microfinance Outputs & Udayan Branch Office System Exports & Evaluated CO-wise Changes Report (member balances, OTR, PAR, write-offs); audited POMIS-4 and POMIS-5 statements. \\ \hline
13--17 Sep 2026 & C-RANGPUR Baseline Data Collection Rotation & Moushumi Akter (MDO, C-RANGPU Project, 01745761900) & Observed field enumeration; tested KoBo Enketo web application; reviewed respondent consent protocols. \\ \hline
\rowcolor{RDRSLight} 20--24 Sep 2026 & Branch Office Visits \& Field Verification & Maydul Islam (01730328212); Md.~Ziaur Rahman (01730328234) & Inspected passbooks against ledgers; audited loan disbursement receipts; followed up overdue recovery entries. \\ \hline
27 Sep 2026 & Human Resource Management Rotation & Md.~Ziaur Rahman (Sr.~Coordinator-HR, 01730328011) & Reviewed HR database systems, staff performance metrics, biometric attendance logs, and field leave tracking. \\ \hline
\rowcolor{RDRSLight} 28 Sep 2026 & Agriculture Dept.\ \& Rupantar Project Briefing & Dr.~Salah Uddin (01716263216); Pallab Goswami (01738 375951) & Analyzed Climate Smart Agriculture (CSA) monitoring; reviewed homestead gardening and flood-resilient crop data. \\ \hline
29 Sep 2026 & Social Development \& Federation Review & Rashedul Arefeen (Team Leader SD, 01730328044) & Mapped governance of 406 Union Federations; examined civic claims tracking and local dispute settlement logs. \\ \hline
\rowcolor{RDRSLight} 30 Sep 2026 & Safeguarding Policy \& Incident Workflow & Rumana Azim Ansari (Coordinator, 01730328039) & Reviewed safeguarding complaint registers, whistleblower protection mechanisms, and confidential reporting paths. \\ \hline
01 Oct 2026 & Admin \& General Service Dept.\ Rotation & Maj Shammita Kabir (Retd) (Head of Admin \& GS) & Inspected procurement registers, warehouse inventory records, vehicle logbooks, and fixed-asset tagging files. \\ \hline
\rowcolor{RDRSLight} 04--16 Oct 2026 & Intensive MIS Practice: Survey Design \& HCVI & Dabasish Mahata; Md.~Ajazul Haque (Manager, MEL) & Conducted KII on C-RANGPUR pipeline; analyzed HCVI formula weighting (30/30/40); mapped F01-112 indicator matrix. \\ \hline
19--26 Oct 2026 & Access to Justice (A2J) Project Walkthrough & Project Focal; DPO; Project Paralegals & Evaluated formal/informal justice split; audited April 2025 Rangpur court docket count (77,849 cases across Ejlash/Serestha). \\ \hline
\rowcolor{RDRSLight} 28 Oct -- 05 Nov 2026 & Archival Audit: Somriddhi Monitoring Report & MEL Department Internal Archives & Documented last-mile data gaps: 89\% register vs.\ 100\% MIS attendance; 11/74 teacher turnover; 65\% missing grant vouchers. \\ \hline
08--12 Nov 2026 & Risk \& Compliance Systems Walkthrough & Microfinance Operations \& Compliance Staff & Evaluated Credit Shield Insurance workflow with Guardian Life; audited AML/CFT database parameters under BFIU rules. \\ \hline
\rowcolor{RDRSLight} 15--25 Nov 2026 & Report Synthesis \& Learning Consultation & Mr.~Mamunur Rashid (Head of MEL) & Synthesized field findings; drafted strategic recommendations; presented draft report for departmental review. \\ \hline
26 Nov 2026 & Formal Internship Presentation Session & MEL Department Leadership \& Staff & Delivered formal presentation on MIS evaluation findings at Rangpur Central Coordination Office. \\ \hline
\rowcolor{RDRSLight} 30 Nov 2026 & Final Report Submission & Academic Supervisor \& RDRS Management & Formally submitted final internship report to BRUR faculty and RDRS MEL leadership. \\ \hline
\end{longtable}
}

% ================= CHAPTER 5 =================
\chapter{Recommendations}

\chapterglance{Actionable Strategic Interventions}{Three prioritized, operational improvements for RDRS monitoring: (1) mandatory register reconciliation sign-off, (2) unified KoBo single-capture pipeline, and (3) institutional knowledge continuity.}

The following recommendations arise directly from the empirical gaps identified in Chapter~\ref{chap:analysis}. They are prioritized by operational urgency.

\section{Recommendation 1: Institute a Mandatory Register-to-MIS Reconciliation Sign-off}
\textbf{The Problem:} The most critical data gap observed at RDRS is the divergence between physical field registers and digital monthly reports (such as the 89\% register attendance entered as 100\% in central MIS).

\textbf{Actionable Solution:}
\begin{enumerate}
\item Introduce a standardized, one-page ``Data Reconciliation Sign-Off Sheet'' into the monthly branch monitoring routine.
\item Before any branch manager or credit organizer can digitally close their monthly G~Banker or project workbook, an independent Monitoring Officer or Area Manager must physically sample at least 10\% of primary field registers (savings passbooks, school roll books, training registers) and compare them line-by-line against digital entries.
\item Both the recording officer and the reviewing monitoring officer must physically sign and timestamp this reconciliation certificate. If a discrepancy exceeds 5\%, the digital record cannot be locked.
\end{enumerate}

\section{Recommendation 2: Automate an End-to-End KoBo-to-Power~BI Pipeline}
\textbf{The Problem:} In development projects like C-RANGPUR, survey data is exported from KoBo into intermediate Excel sheets where formulas and weights are manually copied before being loaded into Power~BI. This creates multiple opportunities for human calculation error.

\textbf{Actionable Solution:}
\begin{enumerate}
\item Configure direct API webhooks from KoBo Toolbox directly into Microsoft Power~BI or an intermediate cloud database (such as PostgreSQL).
\item Embed the Household Climate Vulnerability Index formula ($\mathrm{HCVI} = 0.30\mathrm{EI} + 0.30\mathrm{SI} + 0.40[100-\mathrm{ACI}]$) and indicator calculation logic directly into the automated Power~BI transformation pipeline (using Power Query/DAX).
\item Establish an automated data-cleaning validation script that automatically flags incomplete responses, contradictory inputs, and codebook anomalies (such as inconsistent project naming) upon submission.
\end{enumerate}

\section{Recommendation 3: Codify Institutional Knowledge and Enforce Dual Ownership}
\textbf{The Problem:} While G~Banker has minimized the extreme key-person dependency of the 1990s Micro Bank era, knowledge of complex indicator algorithms and custom queries remains concentrated among a few senior officers.

\textbf{Actionable Solution:}
\begin{enumerate}
\item Develop and publish formal Standard Operating Procedures (SOPs) defining every indicator, ratio numerator, denominator, and calculation rule across all five program areas.
\item Implement a mandatory ``Dual-Ownership Policy'' for every recurring MIS report family, ensuring that at least two trained staff members possess full operational competency to extract, verify, and publish those metrics.
\item Digitize the Credit Shield Insurance dossier transmission, replacing slow physical paper mailing with secure scanned PDF uploads in G~Banker to compress settlement times from weeks to days.
\end{enumerate}

\keybox{\small\textbf{\color{RDRSTeal}Strategic Priority:} If organizational resources permit only one intervention, RDRS should prioritize \textbf{Recommendation 1}. Implementing a physical register-to-system reconciliation sign-off directly eliminates the largest data-integrity risk observed in the field.}

% ================= CHAPTER 6 =================
\chapter{Conclusion}

\chapterglance{Concluding Synthesis}{Reflections on information systems in rural development: why human verification, supervisory routines, and institutional discipline matter more than software brand names.}

My three months inside the MEL Department of RDRS Bangladesh provided a compelling education in how management information systems operate in the development sector. I entered the placement expecting to evaluate software code, database schemas, and server configurations. I emerged realizing that in rural development, MIS is fundamentally about organizational plumbing and human discipline.

The most sophisticated cloud platforms and mobile applications are useless if a field worker fails to verify a passbook, or if an enumerator skips a household in a flood-prone char. RDRS has demonstrated impressive technical progress: its transition from the single-developer fragility of Micro Bank to the centralized automation of G~Banker is a textbook example of successful system evolution. Today, microfinance managers can evaluate credit performance across 288 branches with unprecedented speed and consistency.

The remaining frontier lies at the boundary where software meets the field. Closing the gap between what physical registers say and what central dashboards report requires institutional safeguards: mandatory reconciliation sign-offs, automated data pipelines, and clear procedural ownership. By reinforcing these supervisory checks, RDRS Bangladesh can ensure that its management decisions remain grounded in the true, unvarnished realities of the communities it serves.

% ================= REFERENCES =================
\begin{thebibliography}{99}
\setlength{\itemsep}{1.5pt}\setlength{\parskip}{0pt}
\bibitem{rdrs22} RDRS Bangladesh, \textit{Annual Report 2021--22: Fifty Years of Building Resilience}, Rangpur and Dhaka, 2022.
\bibitem{rdrsstrat} RDRS Bangladesh, \textit{6th Strategic Plan (2023--2027): Sustainable Development through Empowerment}, Dhaka, 2023.
\bibitem{mel2026} RDRS MEL Department, \textit{Departmental Field Monitoring Schedule: September 2026}, Central Coordination Office, Rangpur, 2026.
\bibitem{som23} RDRS Bangladesh, \textit{Field Monitoring Report: Somriddhi Programme (PKSF Funded)}, Panchagarh and Kurigram Field Audits, May 2023.
\bibitem{uday} RDRS Bangladesh, \textit{Microfinance Program Operations Manual: Guidelines for Udayan, Unnoti, Water Credit, and MSDS Savings}, Microfinance Department, 2024.
\bibitem{shield} RDRS Bangladesh and Guardian Life Insurance Limited, \textit{Operational Agreement and Underwriting Guidelines: Credit Shield Microinsurance Policy}, July 2024.
\bibitem{aml} RDRS Bangladesh, \textit{Anti-Money Laundering (AML) and Combating the Financing of Terrorism (CFT) Policy Manual}, effective January 2022.
\bibitem{laudon} K.~C.~Laudon and J.~P.~Laudon, \textit{Management Information Systems: Managing the Digital Firm}, 16th ed., Pearson Education, 2020.
\bibitem{oecd} OECD-DAC, \textit{Glossary of Key Terms in Evaluation and Results Based Management}, Organisation for Economic Co-operation and Development, Paris, 2010.
\bibitem{crangq} RDRS MEL Department, \textit{Baseline Questionnaires of C-RANGPUR: Indicator Result Measurement}, KoBo Toolbox Enketo Application, September 2026.
\bibitem{crangm} RDRS Bangladesh, \textit{F01-112 Enclosure Application Monitoring Matrix: C-RANGPUR Project}, Monitoring Planning Spreadsheet, August 2026.
\bibitem{crangkii} Md.~Ajazul Haque, Manager (MEL), \textit{Technical Key Informant Interview on C-RANGPUR Survey Design, HCVI Formulation, and KoBo Pipeline}, Rangpur, October 2026.
\bibitem{a2jdoc} RDRS Bangladesh and German GIZ, \textit{Access to Justice for Women (A2J for Women): Legal Assistance, Court Audits, and Case Management Deck}, Rangpur, April 2025.
\end{thebibliography}

% ================= APPENDICES =================
\appendix
\chapter{Appendices}

\section{Official Internship Program Schedule at RDRS Bangladesh}\label{sec:official-schedule}
The internship followed a structured three-month placement plan prepared and sanctioned by the Monitoring, Evaluation, and Learning (MEL) Department at RDRS Bangladesh Central Coordination Office, Rangpur. The complete program is reproduced below.

\vspace{0.2cm}
\highlightbox{
\small
\textbf{\color{RDRSBlue}Schedule for the Internship Program at RDRS Bangladesh}\\[2pt]
\textbf{Participant:} Student of BBA, Department of Management Information Systems, Begum Rokeya University, Rangpur\\
\textbf{Duration:} 01 September 2026 to 30 November 2026\\[4pt]
\textbf{Purpose:} The purpose of this internship program is to provide practical exposure to MIS practices in project monitoring for RDRS Bangladesh. The program aims to bridge the gap between theoretical knowledge and real-world practice.
}
\vspace{0.2cm}

{\small\setlength{\tabcolsep}{4pt}
\begin{longtable}{|>{\raggedright\arraybackslash}p{1.9cm}|>{\raggedright\arraybackslash}p{2.3cm}|>{\raggedright\arraybackslash}p{3.8cm}|>{\raggedright\arraybackslash}p{3.5cm}|>{\raggedright\arraybackslash}p{2.2cm}|}
\caption{Official RDRS Internship Program Schedule (01 Sept -- 30 Nov 2026)}\label{tab:official-schedule}\\
\hline
\rowcolor{RDRSBlue}\headercell{Date} & \headercell{Location} & \headercell{Main Activities} & \headercell{Responsible / Contact Person} & \headercell{Remarks / Output} \\ \hline
\endfirsthead
\hline
\rowcolor{RDRSBlue}\headercell{Date} & \headercell{Location} & \headercell{Main Activities} & \headercell{Responsible / Contact Person} & \headercell{Remarks / Output} \\ \hline
\endhead
\hline
\endfoot
\hline
\endlastfoot
01 September 2026 & Central Coordination Office, Rangpur & Orientation and Reporting Guidelines & Mr.~Mamunur Rashid, Head of MEL \newline Dabasish Mahata, Coordinator-Monitoring & Induction completed; placement scope mapped \\ \hline
\rowcolor{RDRSLight} 02--10 September 2026 & Central Coordination Office, Rangpur & Overview of RDRS Bangladesh Programs and Projects & Mr.~Iftekharul Islam, Manager-Monitoring \newline Fahmida Mostafa, Senior MEL Officer & Five thematic areas and MIS toolchains reviewed \\ \hline
13--17 September 2026 & C-RANGPU Project & Baseline Data Collection of C-RANGPU Project & Moushumi Akter, MDO, CRANGPU Project \newline Phone: 01745761900 & Field questionnaire and KoBo setup inspected \\ \hline
\rowcolor{RDRSLight} 20--24 September 2026 & Rangpur Region & Branch Office Visit, Observe Data Verification, and Documentation Process & Maydul Islam, Senior Monitoring Officer (01730328212) \newline Md.~Ziaur Rahman, Senior Officer, Admin \& MIS (01730328234) & Passbook checks; manual-to-digital loan reconciliation \\ \hline
27 September 2026 & HR Department, Rangpur & Know about HRM Policy and Process & Md.~Ziaur Rahman, Sr.~Coordinator-HR \newline Phone: 01730328011 & HR database, attendance, and staffing files \\ \hline
\rowcolor{RDRSLight} 28 September 2026 & Rupantar Project, Rangpur & Activities of Agriculture Department and Rupantar Project & Dr.~Salah Uddin, Team Leader, Environment \& Climate Change (01716263216) \newline Pallab Goswami, Research Associate (01738 375951) & Climate-smart agriculture and homestead monitoring \\ \hline
29 September 2026 & Social Development, Rangpur & Briefing on Social Empowerment Activities and the Federation Concept & Rashedul Arefeen, Team Leader (SD) \newline Phone: 01730328044 & 406 Union Federations governance reviewed \\ \hline
\rowcolor{RDRSLight} 30 September 2026 & Safeguarding Department, Rangpur & RDRS Safeguarding Policy and Practice & Rumana Azim Ansari, Coordinator-Safeguarding \newline Phone: 01730328039 & Safeguarding registers and incident tracking \\ \hline
01 October 2026 & Admin and General Service Dept., Rangpur & RDRS Administrative and General Service activities (Procurement, GH, Store \& GS) & Maj Shammita Kabir (Retd), Head of Admin \& General Service & Asset registers, inventory, and procurement records \\ \hline
\rowcolor{RDRSLight} 04 October -- 12 November 2026 & Rangpur Central Coordination Office & MIS practice in MEL Department & Dabasish Mahata, Coordinator-Monitoring \newline Ajazul Haque, MEL Department & Deep practice: KoBo forms, HCVI scoring, Power~BI, G~Banker, and A2J court data \\ \hline
15--25 November 2026 & Rangpur Central Coordination Office & Report preparation and sharing of learning with RDRS Bangladesh & Mamunur Rashid, Head of MEL & Draft report compiled and reviewed with MEL team \\ \hline
\rowcolor{RDRSLight} 26 November 2026 & Rangpur Central Coordination Office & Internship Presentation Session & Mamunur Rashid, Head of MEL & Presentation delivered to MEL officers and management \\ \hline
30 November 2026 & Rangpur Central Coordination Office & Final Report Submission & Mamunur Rashid, Head of MEL & Formal submission to RDRS and university supervisor \\ \hline
\end{longtable}
}

\section{Semi-Structured Field Interview Guide}\label{sec:interview-guide}
The following core question checklist guided interviews and system walkthroughs with MEL officers, monitoring coordinators, and branch operations staff:
\begin{enumerate}
\setlength{\itemsep}{1.5pt}\setlength{\parskip}{0pt}
\item \textbf{Data Capture Tools:} Which specific instruments (paper registers, Excel templates, mobile apps) are used to collect primary monitoring data in this project?
\item \textbf{Verification and Validation:} Who inspects and cross-checks raw field data before it is aggregated? What documentary evidence is required for sign-off?
\item \textbf{Information Flow:} Through which administrative tiers (Branch $\to$ Area $\to$ Region $\to$ Central Coordination Office $\to$ Executive Leadership) does monitoring information travel, and in what format?
\item \textbf{Software Architecture and Automation:} What functions are performed automatically by central software platforms (e.g., G~Banker, Power~BI) versus manual spreadsheet calculations?
\item \textbf{Data Integrity Bottlenecks:} What are the most common discrepancies observed between field registers and summary MIS reports? Where do reporting delays typically originate?
\end{enumerate}

\section{Glossary of Operational, Legal, and Bangla Terms}\label{sec:glossary}
\begin{itemize}
\setlength{\itemsep}{1.5pt}\setlength{\parskip}{0pt}
\item \textbf{Aposh-joggo:} Offenses that are legally compoundable under the Bangladesh Code of Criminal Procedure (CrPC), allowing lawful settlement through mutual compromise, alternative dispute resolution (ADR), or court-annexed mediation.
\item \textbf{Ejlash:} The open courtroom and elevated bench where formal judicial hearings and trials take place before a judge or magistrate.
\item \textbf{Federation:} Union-level community-owned apex organizations established by RDRS, composed of grassroots village group representatives to manage local advocacy and dispute settlement.
\item \textbf{Monga:} The traditional seasonal lean period in northwestern Bangladesh between crop planting and harvest (typically September to November), historically characterized by severe rural unemployment and food insecurity.
\item \textbf{Samity:} A grassroots village microcredit group consisting of 15--30 female members who meet weekly for savings deposits and joint-liability loan repayments.
\item \textbf{Serestha:} The court clerk's administrative office where physical court records, case dockets, cause lists, and historical case registers are archived.
\end{itemize}

\section{C-RANGPUR Verification Notes and Field Follow-up}\label{sec:crangpur-notes}
The analysis of the C-RANGPUR baseline in Chapter~\ref{chap:analysis} was based on the MEL questionnaire file, the F01-112 matrix spreadsheet, and walkthroughs with Md.~Ajazul Haque. The following technical items are documented for ongoing operational verification:
\begin{enumerate}
\setlength{\itemsep}{1.5pt}\setlength{\parskip}{0pt}
\item \textbf{Sampling Coverage:} Official project documentation should confirm final Upazila and Union breakdown across the six target districts (Rangpur, Lalmonirhat, Kurigram, Sirajganj, Gaibandha, and Jamalpur).
\item \textbf{Reporting Frequency Reconciliation:} In matrix file F01-112, Indicator 1.1 lists both ``semi-annual'' and ``annual'' reporting frequencies; future revisions should clarify whether semi-annual data feeds internal dashboards while annual metrics are reserved for donor reporting.
\item \textbf{Automated Index Calculation:} Transitioning HCVI scoring from manual Excel workbooks into an automated script within the KoBo--Power~BI pipeline will prevent formula transcription errors.
\end{enumerate}

\section{Access to Justice (A2J) Field Notes}\label{sec:a2j-notes}
Data on the A2J for Women project was gathered from the project presentation deck, case audit records, and court walkthroughs. The following items are maintained on record:
\begin{enumerate}
\setlength{\itemsep}{1.5pt}\setlength{\parskip}{0pt}
\item \textbf{Staffing Allocation:} The project employs 6 paralegals, 1 District Project Officer, 1 Finance/Admin Officer, 3 support staff, and 1 Project Focal (discussed headcount of 12 against a slide notation of 13, pending final HR master list confirmation).
\item \textbf{Mediation Referrals:} Track court referrals to panel mediators approved under the Senior District and Sessions Judge's Office Order of 18 March 2025 to evaluate the caseload reduction impact on the formal court docket.
\end{enumerate}

\end{document}
